%%% Hlavní soubor. Zde se definují základní parametry a odkazuje se na ostatní části. %%%

%% Verze pro jednostranný tisk:
% Okraje: levý 40mm, pravý 25mm, horní a dolní 25mm
% (ale pozor, LaTeX si sám přidává 1in)
\documentclass[12pt,a4paper]{report}
\setlength\textwidth{145mm}
\setlength\textheight{247mm}
\setlength\oddsidemargin{15mm}
\setlength\evensidemargin{15mm}
\setlength\topmargin{0mm}
\setlength\headsep{0mm}
\setlength\headheight{0mm}
% \openright zařídí, aby následující text začínal na pravé straně knihy
\let\openright=\clearpage

%% Pokud tiskneme oboustranně:
% \documentclass[12pt,a4paper,twoside,openright]{report}
% \setlength\textwidth{145mm}
% \setlength\textheight{247mm}
% \setlength\oddsidemargin{15mm}
% \setlength\evensidemargin{0mm}
% \setlength\topmargin{0mm}
% \setlength\headsep{0mm}
% \setlength\headheight{0mm}
% \let\openright=\cleardoublepage

%% Pokud používáte csLaTeX (doporučeno):
%\usepackage{czech}
%% Pokud nikoliv:
\usepackage[english,czech]{babel}
\usepackage[T1]{fontenc}

%% Použité kódování znaků: obvykle latin2, cp1250 nebo utf8:
\usepackage[utf8]{inputenc}

%% Ostatní balíčky
\usepackage{graphicx}
\usepackage[bottom]{footmisc}
\usepackage{amsthm,amsmath,amssymb}

%% Balíček hyperref, kterým jdou vyrábět klikací odkazy v PDF,
%% ale hlavně ho používáme k uložení metadat do PDF (včetně obsahu).
%% POZOR, nezapomeňte vyplnit jméno práce a autora.
\usepackage[unicode]{hyperref}   % Musí být za všemi ostatními balíčky
\hypersetup{pdftitle=Využitie genetických algoritmov pri tvorbe rozhodovacích stromov}
\hypersetup{pdfauthor=Lukáš Šurín}
\usepackage[all]{hypcap}

%%% Drobné úpravy stylu

% Tato makra přesvědčují mírně ošklivým trikem LaTeX, aby hlavičky kapitol
% sázel příčetněji a nevynechával nad nimi spoustu místa. Směle ignorujte.
\makeatletter
\def\@makechapterhead#1{
  {\parindent \z@ \raggedright \normalfont
   \Huge\bfseries \thechapter. #1
   \par\nobreak
   \vskip 20\p@
}}
\def\@makeschapterhead#1{
  {\parindent \z@ \raggedright \normalfont
   \Huge\bfseries #1
   \par\nobreak
   \vskip 20\p@
}}
\makeatother

% Toto makro definuje kapitolu, která není očíslovaná, ale je uvedena v obsahu.
\def\chapwithtoc#1{
\chapter*{#1}
\addcontentsline{toc}{chapter}{#1}
}
\newtheorem{def-sk}{Definícia}
\newtheorem*{remark-sk}{Pozorovanie}

\usepackage{caption}
\usepackage{listings} 
\usepackage{subcaption}
\usepackage{algorithm}
\usepackage[noend]{algpseudocode}
\renewcommand{\algorithmicforall}{\textbf{for each}}
\newcommand{\var}[1]{\mathit{#1}}
\newcommand{\func}[1]{\mathrm{#1}}
\usepackage[
	backend=biber,
	style=iso-numeric,
	babel=other,
	sortlocale=cs_CZ,
	bibencoding=UTF8
	]{
	biblatex
	}
\bibliography{bib/wiki,bib/online,bib/kap1,bib/kap2}
\begin{document}

% Trochu volnější nastavení dělení slov, než je default.
\lefthyphenmin=2
\righthyphenmin=2

%%% Titulní strana práce

\pagestyle{empty}
\begin{center}

\large

Univerzita Karlova v Praze

\medskip

Matematicko-fyzikální fakulta

\vfill

{\bf\Large DIPLOMOVÁ PRÁCE}

\vfill

\centerline{\mbox{\includegraphics[width=60mm]{../img/logo.pdf}}}

\vfill
\vspace{5mm}

{\LARGE Lukáš Šurín}

\vspace{15mm}

% Název práce přesně podle zadání
{\LARGE\bfseries Využitie genetických algoritmov pri tvorbe rozhodovacích stromov}

\vfill

% Název katedry nebo ústavu, kde byla práce oficiálně zadána
% (dle Organizační struktury MFF UK)
Katedra softwaru a výuky informatiky

\vfill

\begin{tabular}{rl}

Vedoucí diplomové práce: & RNDr. František Mráz, CSc. \\
\noalign{\vspace{2mm}}
Studijní program: & Informatika \\
\noalign{\vspace{2mm}}
Studijní obor: & ITI \\
\end{tabular}

\vfill

% Zde doplňte rok
Praha 2015

\end{center}

\newpage

%%% Následuje vevázaný list -- kopie podepsaného "Zadání diplomové práce".
%%% Toto zadání NENÍ součástí elektronické verze práce, nescanovat.

%%% Na tomto místě mohou být napsána případná poděkování (vedoucímu práce,
%%% konzultantovi, tomu, kdo zapůjčil software, literaturu apod.)

\openright

\noindent
Poděkování.

\newpage

%%% Strana s čestným prohlášením k diplomové práci

\vglue 0pt plus 1fill

\noindent
Prohlašuji, že jsem tuto diplomovou práci vypracoval(a) samostatně a výhradně
s~použitím citovaných pramenů, literatury a dalších odborných zdrojů.

\medskip\noindent
Beru na~vědomí, že se na moji práci vztahují práva a povinnosti vyplývající
ze zákona č. 121/2000 Sb., autorského zákona v~platném znění, zejména skutečnost,
že Univerzita Karlova v Praze má právo na~uzavření licenční smlouvy o~užití této
práce jako školního díla podle §60 odst. 1 autorského zákona.

\vspace{10mm}

\hbox{\hbox to 0.5\hsize{%
V ........ dne ............
\hss}\hbox to 0.5\hsize{%
Podpis autora
\hss}}

\vspace{20mm}
\newpage

%%% Povinná informační strana diplomové práce

\vbox to 0.5\vsize{
\setlength\parindent{0mm}
\setlength\parskip{5mm}

Názov práce:
Využitie genetických algoritmov pri tvorbe rozhodovacích stromov
% přesně dle zadání

Autor:
Lukáš Šurín

Katedra:
Katedra softwaru a výuky informatiky

Vedoucí diplomové práce:
RNDr. František Mráz, CSc.

Abstrakt:
Rozhodovacie stromy sú úznávanou a široko používanou technikou pri spracovaní a analyzovaní dát. Tieto stromy sú konštruované typickými a všeobecne známymi indukčnými technikami (ID3, C4.5, C5.0, CART, CHAID, MARS, ...). Úspešnosť takto vytvorených stromov nie je vždy dokonalá a často poskytuje určitú rezervu. Pri obtiažnejších kritériách je dokonca indukcia skoro nemožná. V tejto práci sa budeme zaoberať práve týmito rozhodovacími stromami, konkrétne ich vytváraním. Spomínanej rezervy sa snažíme zbaviť pomocou metaheuristiky, genetických algoritmov, ktorá sa využíva pri rôznych optimalizačných úlohách. Práca taktiež obsahuje moju navrhnutú implementáciu algoritmu vo forme zásuvného modelu do prostredia Weka. Porovnanie takto vytvorených stromov so známymi algoritmami tvorby stromov spolu s dosiahnutými výsledkami sú neoddeliteľnou súčasťou tejto práce.

Klíčová slova:
rozhodovacie stromy, genetické algoritmy, indukcia stromov, weka, strojové učenie 

\vss}\nobreak\vbox to 0.49\vsize{
\setlength\parindent{0mm}
\setlength\parskip{5mm}

Title:
Applying genetic algorithms for decision trees induction

Author:
Lukáš Šurín

Department:
Department of Software and Computer Science Education

Supervisor:
RNDr. František Mráz, CSc.

Abstract:
Decision trees are recognized and widely used technique for processing and analyzing data. These trees are designed with typical and generally known inductive techniques (ID3, C4.5, C5.0, CART, CHAID, MARS, ...). Predictive power of created trees is not always perfect and often provides a little reserve. Induction of trees with difficult criterias is even harder and sometime impossible. In this paper we will deal with decision trees, namely their creation. We get rid of mentioned reserve by metaheuristic, genetic algorithms, which is used in all types of optimalization. The work also includes my proposed implementation of the algorithm in the form of plug-in into Weka environment. Comparison of generated trees with known algorithms together with the achieved results are an integral part of this work.

Keywords:
decision trees, genetic algorithms, tree induction, weka, machine learning

\vss}

\newpage

%%% Strana s automaticky generovaným obsahem diplomové práce. U matematických
%%% prací je přípustné, aby seznam tabulek a zkratek, existují-li, byl umístěn
%%% na začátku práce, místo na jejím konci.

\openright
\pagestyle{plain}
\setcounter{page}{1}
\tableofcontents

%%% Jednotlivé kapitoly práce jsou pro přehlednost uloženy v samostatných souborech

\renewcommand{\figurename}{Obrázok}
\chapter{Úvod}
\setlength{\parskip}{0.5em}
\setlength{\parindent}{0cm}
Spracovanie dát je v poslednej dobe nutnou súčasťou množstva procesov (priemyselných, vedeckých, lekárskych a pod). So zvyšovaním množstva, tak ako aj zložitosti dát, je nutné zaručiť určitú mieru automatického spracovania. Hlavnou úlohou je získať čo najrelevantnejšie informácie, pre danú oblasť, pričom ľudský faktor by mal zohrávať minimálnu rolu. Obor, ktorý rieši takéto úlohy sa nazýva Dobývanie znalostí\footnote{Dobývanie znalosti je len jednou zo súčastí oveľa bohatšieho, väčšieho procesu nazývaného proces dobývania znalostí z databází, anglicky knowledge discovery in databases (KDD)}. LinkedIn ako aj Facebook sú príkladmi, kde je táto technika využívaná na automatické odporúčania nových priateľov, ktorých ešte nemáme v kontaktoch.

Celý proces spracovania dát sa zlepšovaním výpočtovej techniky, ktorú si je ľahké a lacné zadovážiť i s prípadnou paralelizáciou, radikálne zlepšil. Došlo k zrýchleniu tvorby modelov, čo v konečnom dôsledku zlacnilo celý tento postup. Zlepšenie modelu je ale taktiež dôležitou súčasťou. Vhodná optimalizácia pomocou genetických algoritmov, či iných heuristík môže významne zlepšiť počiatočný stav. Reťazenie, kombinovanie a využívanie známych techník ako bagging, boosting a ďalších je taktiež dôležitou taktikou, ktorá zostáva prínosom pre túto sféru. Dostupné metódy sa ale neustále vylepšujú a to hlavne kvôli ustavične rastúcej zložitosti.

Rozhodovacie stromy, ako aj ďalšie iné techniky (asociačné pravidlá, neurónové siete a pod) sú súčasťou dobývania znalosti. Rozhodovací strom patrí pod skupinu prediktívnych modelov, ktoré môžeme použiť na klasifikáciu\footnote{Klasifikačný strom}, ako aj na regresiu\footnote{Regresný strom}. V praxi je obľúbený hlavne kvôli jeho jednoduchosti a zrozumiteľnosti, ktorá vychádza z jeho priamočiarej stromovej štruktúry. Extrakcia pravidiel je v tomto prípade len otázka prechodu touto štruktúrou. Tieto pravidlá sú ľahko zrozumiteľné a je ich následne možné použiť na klasifikáciu. Väčšina techník z dobývania znalostí je založená na induktívnom učení. V takomto prípade je model skonštruovaný na trénovacích dátach, tak aby bol čo najviac generalizovaný aj pre reálne dáta. Takýto prístup ale neposkytuje vždy dostačujúce prediktívne schopnosti, ktoré by zaručili využiteľnosť tohoto modelu. Veľkým omedzením týchto techník je neexistujúca schopnosť využiť komplikovanejšie kritéria alebo nejakú ich kombináciu.

Keď neberieme do úvahy problém s komplikovanými kritériami, tak je na zlepšenie modelu možné využiť viacero prístupov. Jednoduché zlepšenie schopností modelu sa dá dosiahnuť v poskytnutím zväčšenej trénovacej množiny. Najväčším problémom býva to, že sa pri úlohách stretávame skorej s menším množstvom dát. Bez použitia nejakých ďalších techník, ktoré dokážu vytvárať využiteľné modely aj na týchto omedzených množinách, je tento prístup nepoužiteľný. Tiež treba brať do úvahy aj to, že indukčný algoritmus nemusí vedieť efektívne pracovať s takouto zväčšenou množinou a môže dojsť k preučeniu. Potenciálnym riešením pri stromoch býva v takomto prípade orezanie.

Ďalšou možnosťou je využiť nejaký iný indukčný algoritmus, ktorý by vedel využiť všetky informácie v dátach čo najlepšie. Vytvorenie nového rýchleho a správneho algoritmu ale nie je tou najľahšou úlohou. Existujúce algoritmy sú v tejto dobe na vysokej úrovni a získavajú maximum z týchto dát bez toho, aby sme znížili generalizačné schopnosti.

Veľmi efektívnym riešením by bolo vyberanie len takých dát z trénovacej množiny, ktoré sú pre naše prediktívne schopnosti modelu dôležité. Hlavným problémom je však to, že dopredu nevieme určiť, ktoré dáta budú dôležité.

Na zlepšenie schopností modelu budeme v tejto práci využívať už spomínané genetické algoritmy. Najvýznamnejšou výhodou tohoto prístupu je popri robustnosti, pridanie schopnosti vytvárať také modely, ktoré dokážu využívať aj rôzne komp\-likované kritéria. Nevýhodou je zvyšovanie výpočtovej náročnosti (najmä doby výpočtu), ktorá narastá priamo úmerne so zložitosťou kritérií.

Existuje mnoho článkov, ktoré sa zaoberajú vytváraním rozhodovacích stromov pomocou genetických algoritmov. Ich výsledky sú pomerne zaujímavé. Voľne dostupných, implementovaných riešení je ale málo. Týmto sa teda tento prístup stáva obtiažnejším na použitie.
\section{Cieľ práce}
Táto práca má za úlohu vytvoriť algoritmus na budovanie rozhodovacích stromov využívajúci genetické algoritmy. Ťažiskom práce nebude generovanie stromov kompletne len genetickými algoritmami, ale hlavne využitie dostupných indukčných techník, ktoré nám dajú dobrý základ počiatočnej populácie. Následne budeme genetické algoritmy využívať na dooptimalizovanie stromov a to aj na základe ľubovoľných kritérií alebo ich kombinácii. Takto sa snažíme čo najviac využiť informácie v dátach a zvýšiť prediktívne schopnosti generovaných rozhodovacích stromov. Zároveň sa snažíme tento algoritmus sprístupniť už z nejakého existujúceho voľne dostupného prostredia, ktoré je používané širokou verejnosťou a je naviac nezávislé na platforme. Vytvorený program by mal taktiež zaručovať určitú mieru rozšíriteľnosti. Týmto rozhodnutím bude riešenie dostupné aj pre bežných užívateľov, ktorí budú môcť využívat vyvinuté postupy aj na bežnom osobnom počítači. 

\section{Štruktúra práce}
V kapitole \ref{kap1:DT} najprv uvedieme dobývanie znalosti a predstavíme rozhodovacie stromy. V kapitole \ref{kap2:GA} zavedieme metaheuristiku nazývanú genetické algoritmy. V kapitole \ref{kap3:DTGA} zase pospájame doposiaľ definované pojmy a zostavíme nový genetický algoritmus na vytváranie rozhodovacích stromov. Následne v kapitole \ref{kap4:Implementation} aplikáciu \verb|GenDTLib| i s podstatnými detailami implementácie nášho genetického algoritmu. V kapitole \ref{kap5:Tests} popíšeme, ako sme algoritmus testovali a aké výsledky táto technika produkovala. V tejto kapitole zároveň porovnáme náš genetický algoritmus na tvorbu rozhodovacích stromov so známym indukčným algoritmom C4.5. Ďalej v kapitole \ref{kap6:SimilarWorks} uvedieme popis existujúcich prác, ktoré vytvárajú rozhodovacie stromy genetickými algoritmami. Nakoniec v závere práce \ref{kap:fin} zhodnotíme dosiahnuté výsledky a popíšeme možné budúce rozšírenia práce.

Práca obsahuje viacero dodatkov. V dodatku \ref{kapI} je obsah priloženého média. Dodatok \ref{kapII} slúži ako užívateľská príručka aplikácie \verb|GenDTLib| s popisom parametrov genetického algoritmu, ktoré môžeme zmeniť v konfiguračnom súbore. Nakoniec v dodatku \ref{kapIII} popíšeme návod pre užívateľov, ktorí by chceli rozširovať aplikáciu o nové komponenty implementovaním daných rozhraní.
\newcommand{\infogain}{\mbox{Info\_zisk}}
\newcommand{\entropy}{\mbox{Entropia}}
\newcommand{\ginindex}{\mbox{Gini\_zisk}}
\newcommand{\gini}{\mbox{Gini}}
\newcommand{\gainratio}{\mbox{Koef\_zisku}}
\newcommand{\twoing}{\mbox{Twoing}}

\chapter{Úvod k rozhodovacím stromom}
V tejto kapitole zadefinujeme základné pojmy a koncepty, ktoré budú obsiahnuté v ďalších častiach tejto práce. V nadchádzajúcich podsekciách vysvetlíme postupne všetky dôležité pojmy. V oddieli \ref{kap1:2.1:Data} pomenovanom Dáta zadefinujemem čo sú to dáta, zavedieme pojem atribút a  vysvetlíme základné pojmy, s ktorými budeme pracovať. Dobývanie znalostí, ktoré predstavuje oddiel \ref{kap1:2.2:DataMining}, popíše tento obor so zaradením rozhodovacích stromov v ňom. V oddieli \ref{kap1:2.3:DT} zavedieme základný model rozhodovacích stromov, pre ktorý bude nutné uviesť pár pojmov z teórie grafov. Oddiel \ref{kap1:2.4:DTTypes} zoznamuje s rôznymi typmi rozhodovacích stromov. V oddieli \ref{kap1:2.5:DTTechniques} zhrnieme techniky nazývané indukčné algoritmy na tvorbu stromov.
Nakoniec v oddieli \ref{kap1:2.6:DTUsage} ukážeme využitie rozhodovacích stromov pri predikcii dát.

\section{Dáta}\label{kap1:2.1:Data}
V oddieli popíšeme čo sú to dáta. Dozvieme sa tie najzákladnejšie pojmy s ktorými sa stretneme pri práci s nimi. Informácie sme čerpali z voľne dostupnej knihy  \cite{kap1-DataMiningAndAnalysis} a online zdrojov \cite{wiki-Data,online-Data}.

V prvej časti \ref{kap1:2.1:2.1.1:DataRepresentation} popíšeme, čo sú to dáta a ako ich môžme reprezentovať. Popíšeme základné pojmy, ako je atribút a príznakový vektor. V časti \ref{kap1:2.1:2.1.2:DataAttributes} rozdelíme atribúty na dva druhy - kategoriálne a numerické. Aj napriek jednoduchosti týchto pojmov je ich nutné zaviesť, pretože ich budeme využívať v neskorších kapitolách.

\subsection{Dáta a ich reprezentácia}\label{kap1:2.1:2.1.1:DataRepresentation}
Dáta sú množiny hodnôt, ktoré predstavujú kusy informácií.
Môžu byť zozbierané, odmerané, analyzované a následne vizualizované.
Zozberané dáta môžme nájsť v relačných, tabuľkových databázach. Tento typ dát nazývame štrukturované. V skutočnosti sa však vo väčšine prípadov stretávame presne s ich opakom, a to s neštrukturovanými dátami.
Podľa \cite{kap1-DataMiningAndAnalysis} si dáta možme predstaviť ako maticu $n \times d$, kde $n$ predstavuje počet riadkov a $d$ počet stĺpcov tejto matice. Riadky sú v tomto prípade popisujú objekty a stĺpce sú tvorené hodnotami danej vlastnosti. Matica je daná takto

\begin{center}
Data = 
$\begin{array}{c | c c c c}
& X_{1} & X_{2} & \ldots & X_{d} \\ \hline
\mathbf{x}_{1} & x_{11} & x_{12} & \ldots & x_{1d} \\
\mathbf{x}_{2} & x_{21} & x_{22} & \ldots & x_{2d} \\
\vdots & \vdots & \vdots & \ddots & \vdots \\
\mathbf{x}_{n} & x_{n1} & x_{n2} & \ldots & x_{nd} \\
\end{array}$
\end{center}

$x_{i}$ predstavuje $i$-tý riadok, ktorý obsahuje $d$ hodnôt.
Jeden riadok matice taktiež nazývame príznakový vektor, entita, objekt, transakcia alebo inštancia.
\begin{center}
$\mathbf{x}_{i} = (x_{i1},x_{i2},\ldots,x_{id})$
\end{center}

$X_{j}$ predstavuje $j$-tý stĺpec, ktorý obsahuje $n$ hodnôt.
Na druhú stranu, stĺpec matice označujeme aj ako príznak, vlastnosť, atribút.
\begin{center}
$X_{j} = (x_{1j},x_{2j},\ldots,x_{nj})$
\end{center}

Pri metódach učenia s učiteľom uvažujeme ešte jeden stĺpec. Jeho hodnoty odpovedajú klasifikácii alebo regresii jednotlivých inštancií. Tieto hodnoty sú z domény množiny označovanej $C$. Kvôli konzistencii budeme tento posledný stĺpec ďalej nazývať ako výstupný atribút.
\begin{align}
C &= (c_{1},c_{2},\ldots,c_{m}) \nonumber
\end{align}

Premennú \textit{n} nazývame veľkosť dát a premennú \textit{d} zase dimenzionalita alebo rozmer dát.

Ako sme už spomínali, štrukturované dáta sú tie, ktoré máme uchované v riadkovo-stĺpcových databázach a teda ich ľahko prevedieme do maticového zápisu. Pre komplexnejšie dátové množiny (neštrukturované), ktoré sa vyskytujú napr. v bioinformatike (DNA, proteínové sekvencie,...), je buď nutné použiť iný zápis alebo využiť techniku extrakcie príznakov. Medzi ďaľšie neštrukturované dáta zaraďujeme obrázky, text, audio ale aj video nahrávky.

\subsection{Atribúty}\label{kap1:2.1:2.1.2:DataAttributes}
Doména atribútu je množina hodnôt, ktoré môže atribút nadobudnúť. Atribúty môžme rozdeliť na 2 typy podľa ich domény:
\begin{itemize}
\item \textbf{Kategoriálne} -- doména týchto atribútov sa skladá z konečnej množiny symbolov. Takýmto atribútom môže byť napríklad pohlavie (muž, žena), rodinný stav (slobodný, ženatý), ale aj komplikovanejšie ako je študijný program (informatika, fyzika, matematika, ...). Tieto atribúty je možné rozdeliť ešte do dvoch skupín:
\begin{itemize}
\item \textbf{Nominálne}, keď hodnoty v doméne nie sú usporiadané. Zmysluplné je iba porovnanie na zhodu. Príkladom môže byť pohlavie.
\item \textbf{Usporiadané}, keď sa hodnoty v doméne dajú porovnávať. Takéto hodnoty môžeme porovnávať nie len na zhodu, ale aj porovnávať (hodnota je väčšia, menšia). Ako príklad môžeme uviesť vzdelanie (základné, stredné, vysokoškolské).
\end{itemize}
\item \textbf{Numerické} sú také, ktorých doména je založená na celých alebo reálnych číslach. Jedným z takýchto atribútov môže byť príjem, zostatok na účte, \ldots. 
\end{itemize}

\section{Dobývanie znalostí}\label{kap1:2.2:DataMining}
V tomto oddieli zhrnieme vedomosti z dobývania znalostí, ktorého základ je dôležitý pri zaradení rozhodovacích stromov do správneho oboru.
Znalosti sme čerpali z kníh, ktoré sú vhodné ako úvod do tohoto oboru \cite{kap1-DataMiningAndAnalysis,kap1-DataMiningForMasses,kap1-DataMiningForTrees,kap1-StatisticLearn}. Pomimo týchto kníh sú informácie získané aj z ďaľších zdrojov ako je Wikipédia \cite{wiki-DataMining} ale aj menej známa HTML stránka od Dr. Saed Sayada \cite{online-DataMining}, ktorá stojí za povšimnutie. Všetky materiály sú voľne dostupné online.
 
V časti \ref{kap1:2.2:2.2.1:KDD} popíšeme proces dobývania znalostí z databází, anglicky knowledge discovery in databases (KDD). Hlavné informácie o samotnom dobývaní znalostí zhrnieme v časti \ref{kap1:2.2:2.2.2:DataMineProcess}. Tu ho rozdelíme do štyroch vrstiev. Každú vrstvu vysvetlíme a pri tom do tejto schémy zaradíme rozhodovacie stromy. V časti \ref{kap1:2.2:2.2.3:Taxonomy} predstavíme taxonómiu metód s nadväznosťou na zavedený 4-vrstvový model. Ďalší oddiel \ref{kap1:2.2:2.2.4:Supervised} vysvetlí, čo sú to metódy učenia s učiteľom, rozlíšime jeho dve obmeny a spomenieme tu aj učenie bez učiteľa.
Nakoniec v oddieli \ref{kap1:2.2:2.2.5:Tools} zhrnieme a uvedieme pár základných nástrojov, ktoré môžeme použiť pri práci v oboru dobývania znalostí. 

\subsection{Dobývanie znalostí z databází}\label{kap1:2.2:2.2.1:KDD}
Tiež známe pod skratkou KDD je komplikovaný proces pozostávajúci z viacerých fáz, ktorý ma za úlohu identifikovať nové, využiteľné vzory v dátach. KDD zahŕňa aj samotné dobývanie znalostí, ktoré je pravdepodobne jeho najvýznamnejšou súčasťou. Kvôli prominentnému postaveniu dobývania znalostí v rámci KDD, ich mnoho odborníkov stotožňuje.

\begin{figure}[h]
\centering
\centerline{\mbox{\includegraphics[width=400pt]{../img/kap1/DM-KDD.pdf}}}
\caption{Štrukturovaný KDD proces (prekreslené a preložené z \cite{kap1-DataMiningForTrees}).}\label{fig:dataMineKDD}
\end{figure}

Jednou z úloh v minulosti bolo formalizovať a štandardizovať prístup k dobývaniu znalostí. Na obrázku \ref{fig:dataMineKDD} sú jednotlivé časti sedem-stupňového modelu. Toto rozdelenie ale nie je pevné. Mnoho iných prác a spoločností navrhlo svoje obmeny. Veľké korporácie ako Daimler-Benz, poskytovateľ poistenia OHRA, vývojári softwaru a hardwaru NCR Corp. sa v roku 1999 spojili a vytvorili svoj vlastný model. Týmto vznikol známy CRISP-DM, ktorého štruktúru je možné vidieť na obrázku~\ref{fig:dataMineCRISP} \cite{kap1-DataMiningForMasses}. Ďalším známym modelom je model 5A alebo SEMMA.

\begin{figure}[h]
\centering
\centerline{\mbox{\includegraphics{../img/kap1/DM-CRISP.pdf}}}
\caption{CRISP-DM model (prekreslené a preložené z \cite{kap1-DataMiningForMasses})}\label{fig:dataMineCRISP}
\end{figure}

\subsection{Popis dobývania znalostí}\label{kap1:2.2:2.2.2:DataMineProcess}
Dobývanie znalostí je technika z KDD, ktorá je priamo zodpovedná za vytváranie poznatkov o existujúcich dátach, tak isto ako za predikciu určitých vlastnosti dát. Je to obor, ktorý kombinuje viaceré poznatky zo štatistiky, umelej inteligencie, strojového učenia a čiastočnej znalostí databázových strojov (obrázok \ref{fig:dataMineComb}).

% IMAGE 1
\begin{figure}[h]
\centering
\centerline{\mbox{\includegraphics[width=225pt]{../img/kap1/DM-comb.pdf}}}
\caption{Dobývanie znalostí ako kombinácia rôznych oborov (prekreslené a preložené z \cite{online-DataMining})}\label{fig:dataMineComb}
\end{figure}

Celý proces dobývania znalostí je možné popísať v štyroch vrstvách. Tento model je možné vidieť na obrázku \ref{fig:layerModel}. Každá z nižších vrstiev je dôležitá pre vrstvu vyššie.

% IMAGE 2
\begin{figure}[h]
\centering
\centerline{\mbox{\includegraphics[width=400pt]{../img/kap1/DM-layer.pdf}}}
\caption{Proces dobývania znalostí v 4 vrstvách (prekreslené a preložené z \cite[s. 26]{kap1-DataMiningForTrees})}\label{fig:layerModel}
\end{figure}

Prvá vrstva sa zaoberá už konečnou, výstupnou aplikáciou, pri tvorbe ktorej sme využili všetky predchádzajúce vrstvy. Tieto majú široký záber využitia v podnikateľských sférach. Medzi najvyužívanejšie patrí ohodnocovanie zákazníkov podľa ich príjmu, tak ako aj detekcia fraudov\footnote{podvod alebo nepoctivý trik, ktorým je poškodený zainteresovaný uživateľ}. Na vytvorenie aplikácie použijeme jednu alebo viacero techník, ktoré na našom obrázku predstavujú druhú vrstvu. Táto vrstva sa zaoberá úlohami strojového učenia ako je regresia, zhlukovanie a mnoho iných. Tieto úlohy využívajú rôzne, už konkrétne modely (tretia vrstva), kde patria už spomínané rozhodovacie stromy, neurónové siete, kohonenove mapy, \ldots. Posledná vrstva sa zaoberá tvorbou týchto modelov, v ktorej sa nachádzajú aj indukčné algoritmy tvorby rozhodovacích stromov.


\subsection{Taxonómia metód}\label{kap1:2.2:2.2.3:Taxonomy}
Z predchádzajúceho oddielu vieme, že časti dobývania znalostí sú umiestnené v nejakých vrstvách. V tomto oddieli nás budú zaujímať hlavne vzťahy v rámci druhej a tretej vrstvy. Vertikálne rozlíšenie je uvedené v predchádzajúcom oddieli. Na zistenie nejakých ďalších vzťahov je nutné zaradiť metódy do širších jednotiek ako na obrázku \ref{fig:dataMineParad}. Jeden taký pohľad na rozdelenie metód môžeme nájsť aj v \cite{online-DataMining}. V tomto druhom prípade je taxonómia v konečnom dôsledku funkcionálne rovnaká (niektoré názvy sú iné), ale líšia sa štruktúrou. 

V úlohách je v prvom rade nutné rozlišovať medzi dvoma typmi dobývacích techník, verifikačné a objaviteľské. Každý z nich má svoju metodológiu. Verifikačný sa zaoberá testovaním hypotéz, analýzou rozptylu a ďalšími praktikami. Objaviteľské zase automaticky hľadajú nové pravidlá a vzory.

Verifikačné metódy spoliehajú na určitú predom známu hypotézu, ktorá býva vyhodnotená expertom. Tento prístup nie je až tak spojený s dobývaním znalostí, ako to je pri objaviteľských metódach. Najvýznamnejším dôvodom je ten, že techniky dobývania znalostí by mali vyberať hypotézu z množiny vhodných hypotéz (identifikácia modelu), než aby sa spoliehali na nejakú dopredu známu (odhadnutie modelu).

Medzi takéto objaviteľské techniky patria ako predikčné, tak aj popisné metódy (zhlukovanie, vizualizácia). Popisnými sa hlavne snažíme pochopiť ako a prečo sú dáta umiestnené tak, ako sú nám prezentované. Predikčné zase dokážu predpovedať hodnoty jednotlivých atribútov daného objektu. Indukčné algoritmy sú typickou technikou, ktoré sa používajú pri vytváraní týchto modelov.

\begin{figure}[h]
\centering
\centerline{\mbox{\includegraphics[width=400pt]{../img/kap1/DM-paradigm.pdf}}}
\caption{Taxonómia modelov, prekreslené a preložené z \cite{kap1-DataMiningForTrees})}\label{fig:dataMineParad}
\end{figure}

\subsection{Učenie s učiteľom a bez učiteľa}\label{kap1:2.2:2.2.4:Supervised}
Učenie s učiteľom je oblasťou dobývacích techník, ktoré na vytvorenie modelu potrebuje trénovaciu množinu. Z taxonómie zavedenej v predchádzajúcom oddieli sem patria predikčné metódy. Pri tomto type sa teda snažíme nájsť vzťahy v trénovacích dátach medzi vstupnými atribútmi (nazývanými aj nezávislé premenné) a výstupným, klasifikačným atribútom (tiež nazývaný závislá premenná). Nájdený vzťah v dátach je reprezentovaný určitou štruktúrou a je to v konečnom dôsledku náš model. Takýto model je vytváraný už nejakou konkrétnou metódou. Medzi najzákladnejšie patria rozhodovacie stromy, neurónové siete, SVM, \ldots. Poskytnutím nových inštancií môžeme model využiť na klasifikáciu výstupného atribútu z tých vstupných.
Využitie je veľmi bohaté v rôznych sférach (financie, obchod, služby, \ldots). 

Hlavným cieľom učenia je zlepšenie modelu na nejakej úlohe pomocou predom daných dát. K tomuto cieľu je nutné poznať tri jeho komponenty:
\begin{itemize}
\item Úloha $U$, ktorú chceme učením vylepšiť
\item Dáta $D$, ktoré použijeme pri učení
\item Meradlo výkonu $M$, ktoré je použité pri meraní miery zlepšenia.
\end{itemize} 
Pre lepšie pochopenie môžeme použiť známy problém identifikácie emailov do spamu\footnote{Spamová správa - nevyžiadaná správa, ktorú užívateľovi prišla do emailovej schránky.}. Pre túto vyzerajú komponenty takto:
\begin{itemize}
\item Úlohou $U$ je identifikovať nevyžiadané emaily.
\item Dáta $D$ sú v tomto prípade množiny emailov, v ktorých sa nachádzajú správne aj nesprávne emaily.
\item Meradlom výkonu $M$ je percento nevyžiadaných emailov, ktoré boli klasifikované správne a percento správnych emailov, ktoré boli klasifikované nesprávne.
\end{itemize} 

Trénovacie dáta sú v tomto prípade v trochu odlišnejšom dátovom zápise, ako sme definovali v časti \ref{kap1:2.1:Data}. Nová matica by vyzerala ako $TrainData = (Data | Y)$, kde znak $|$ znamená pripojenie stĺpca $Y$ za maticu $Data$. Iným zápisom môže byť aj $TrainData = \{(\mathbf{x}_{k},y_{k})\}_{k=1}^{n}$, kde $y_{k}$ je hodnota klasifikácie/regresie pre vektor $\mathbf{x}_{k}$. Hodnota $y_{k}$ je z domény množiny $C$ 

Pri modeloch učenia s učiteľom je nutné rozlišovať medzi dvoma základnými typmi podľa typu výstupného atribútu:
\begin{itemize}
\item \textbf{Klasifikačné} mapujú vstupný vektor na dopredu známe triedy. Príkladom môže byť predikovanie rizika rakoviny (žiadne, nízke, vysoké) pacienta podľa jeho lekárskej anamnézy.
\item \textbf{Regresné} mapujú vstupný vektor na reálne čísla. Takýto model dokáže napríklad určiť výšku pôžičky, ktorú môže banka poskytnúť z informácii o žiadateľovi.
\end{itemize} 

Učenie bez učiteľa je presným opakom od učenia s učiteľom. Tento typ vytvára model bez trénovacej množiny. Podľa taxonómie sem môžme zase zaradiť popisné metódy (nie vizualizáciu) z objaviteľských techník. 

\subsection{Nástroje pre DM}\label{kap1:2.2:2.2.5:Tools}
Existuje veľké množstvo nástrojov, ktoré boli vytvorené aby zvládali úlohy z dobývania znalostí. Medzi nimi existuje plno profesionálnych nástrojov určených priamo pre takéto úlohy (Weka, RapidMiner). Iné druhy sú vo forme knižníc do stávajúcich jazykov (Python, Matlab, ...). Stavané sú tak, aby dokázali využiť techniky umelej inteligencie, strojového učenia, štatistiky a iné. Veľa z nich je ale platených, niektoré zase komplikované na inštaláciu, konfiguráciu alebo použitie. Pre učenie, ľahšiu využiteľnosť a hlavne pre účel našej práce sú preto nevhodné. 

\begin{figure}[h]
\centering
\centerline{\mbox{\includegraphics[width=300pt]{../img/kap1/DM-Weka.png}}}
\caption{Náhľad na užívateľské prostredie nástroja Weka}\label{fig:dataMineWeka}
\end{figure}

Najvhodnejšími kandidátmi sú:
\begin{itemize}
\item \textbf{RapidMiner} je napísaný v jazyku Java a poskytuje pokročilé analytické techniky pomocou framework-u založeného na šablónach. Výhodou tohto nástroja je aj to, že používateľ nie je nútený vytvárať žiaden kód. Ponú\-kaný je skorej ako služba, než kus lokálneho softwaru. Popri riešení úloh dobývania znalostí dokáže aj predspracovanie, vizualizáciu dát, štatistické modelovanie a vyhodnocovanie. RapidMiner je jedným z najpoužívanejších a najlepších nástrojov na úlohy z dobývania znalostí. Sila nástroja narastá vďaka využiteľnosti schém, modelov a algoritmov z Weky a R-skriptov. Distribuovaný je pod AGPL licenciou a je voľne stiahnuteľný zo stránky SourceForge.
\item \textbf{Weka} je kolekciou algoritmov pre strojové učenie a dobývanie znalostí. Algoritmy môžu byť volané priamo z užívateľského prostredia programu (obrázok \ref{fig:dataMineWeka}) alebo z vlastne vytvorenej aplikácie bežiacej pod Java prostredím. Obsahuje nástroje na predspracovanie, klasifikáciu, regresiu, asociačné pravidlá a vizualizáciu. Taktiež je použiteľná na vytváranie vlastných schém (prebrané z \cite{online-DataMiningWeka}). Aplikáciu si je možné stiahnuť zadarmo z hlavnej stránky pod GNU General Public licenciou. Tento typ licencie je jedna z najdôležitejších výhod Weky oproti RapidMiner-u. Vďaka nej je možná ľubovoľná úprava nástroja podľa našich požiadavkov, tak ako aj pridanie úplne nových algoritmov.
\end{itemize}

Ďalší kandidáti nie sú samostatnými aplikáciami, ale moduly pre dobývanie znalostí do nejakého dynamického jazyka. Tieto možnosti môžeme považovať za pokročilejšie, keďže na ich využitie je nutné spísať vlastný kód vo forme skriptov.
\begin{itemize}
\item \textbf{Matlab} je jazyk vysokej úrovne a interaktívne prostredie používané širokou skupinou vedcov a technikov \cite{online-DataMiningMatlab}. Príjemné prostredie tohoto jazyka a možnosť interaktívneho ladenia svojich programov zjednodušuje prácu. Rozširujúce moduly nazývané toolboxy\footnote{panel nástrojov, pomenovanie modulov do Matlab-u}, majú široké spektrum zamerania. Medzi najvýznamnejšie moduly na tvorbu modelov a dátovú analýzu pat\-ria štatistický modul (Statistics and Machine Learning Toolbox), modul pre neu\-rónové siete (Neural Network Toolbox) a mnoho ďalších. Tieto moduly sa dajú často ovládať z takzvaných pomocných nástrojov\footnote{pomocný nástroj/wizard je skupina po sebe nasledujúcich obrazoviek, ktoré prevedú užívateľa celým procesom tvorby požadovaného programu, bez ďalšej znalosti vnútornej funkcionality}.
Nevýhodou tohoto prístupu je buď nutnosť vlastniť Matlab, ktorý je spoplatnený, alebo mať alternatívny Octave, ktorým zase strácame príjemné užívateľské prostredie s ďalšími výhodami modulov, ako sú spomenuté pomocné nástroje.
\item \textbf{Python} je vysoko úrovňový dynamický jazyk podobne ako to je u Matlab-u \cite{wiki-Python}. Je to obľúbeným jazykom vedcov a jeho hlavnou filozofiou je zaručiť čitateľnosť kódu. Python podporuje mnoho programátorských paradigmát (objektovo orientovaný, imperatívny, funkcionálny, ...). Taktiež k nemu existuje mnoho knižníc na zjednodušenie práce v širokej škále oborov. Najoblúbenejší nástroj pre štatistické a strojové učenie je obsiahnutý v balíčku scikit-learn \cite{kap1-Scikit}. Tento balík v sebe viaže efektívne a jednoduché nástroje pre dátovú analýzu aj dobývanie znalostí. Stavia na obľúbených knižniciach NumPy, SciPy a matplotlib. Jazyk Python je k tomu voľne dostupný a scikit-learn je so svojou BSD licenciou skvelým kandidátom aj pre obchodné účely. 
\end{itemize}


\section{Základné pojmy rozhodovacích stromov}\label{kap1:2.3:DT}
V tomto oddieli vysvetlíme základné pojmy, ktoré sa používajú v spojitosti s rozhodovacími stromami. Najprv zadefinujeme čo je strom a čo znamenajú jednotlivé jeho časti v rozhodovacom strome. Informácií o tejto metóde existuje veľa, no základné pojmy sme hlavne čerpali z kníh \cite{kap1-DataMiningForTrees,kap1-DecisionTree} a \cite[s. 481-498]{kap1-DataMiningAndAnalysis}. Ďalšími použitými zdrojmi bola Wikipédia \cite{wiki-DecisionTree} a scikit-learn stránka \cite{online-DecisionTreeScikit}, na ktorej sa nachádzajú zaujímavé príklady.

V časti \ref{kap1:2.3:2.3.1} popíšeme čo je strom. Rozhodovacie stromy potom na základe grafovej štruktúry stromu zavedieme v časti \ref{kap1:2.3:2.3.2}.

\subsection{Strom}\label{kap1:2.3:2.3.1}
Na bližšie zadefinovanie rozhodovacieho stromu je potrebný aspoň minimálny základ z teórie grafov. Základné pojmy sme čerpali z knihy \cite{kap1-KapitolyDiskretka}.
\begin{def-sk}[Graf]\label{kap1:2.3:2.3.1:graf}
Graf je usporiadaná dvojica (V,E), kde V je nejaká neprázdna množina a E je množina dvojbodových podmnožín množiny V. Prvky množiny V sa nazývajú vrcholy grafu G a prvky množiny E sú hrany grafu G.
\end{def-sk}

\begin{figure}[h]
\centering
\centerline{\mbox{\includegraphics[width=300pt]{../img/kap1/DT-graphs.pdf}}} 
\caption{Príklady neorientovaných grafov. Čierne body predstavujú vrcholy a čiary medzi vrcholmi reprezentujú hrany z $E$.}\label{fig:decisionTreeGraphs}
\end{figure}

\begin{def-sk}[Orientovaný graf]\label{kap1:2.3:2.3.1:orient-graf}
Orientovaný graf G je dvojica (V,E), kde V je nejaká neprázdna množina a E je podmnožina kartézskeho súčinu $V \times V$. Prvky množiny V sa nazývajú vrcholy grafu G a prvky množiny E nazývame orientované hrany. Orientovaná hrana má tvar (x,y). Hovoríme, že orientovaná hrana vychádza z x a končí v y.
\end{def-sk}

Príklady obyčajných(neorientovaných) a orientovaných grafov je možné vidieť na obrázkoch \ref{fig:decisionTreeGraphs} a \ref{fig:decisionTreeOrientGraphs}.

\begin{figure}[h]
\centering
\centerline{\mbox{\includegraphics[width=300pt]{../img/kap1/DT-orientgraphs.pdf}}}
\caption{Príklady orientovaných grafov. Čierne body predstavujú vrcholy a čiary medzi vrcholmi reprezentujú orientované hrany z $E$. Šípka určuje orientáciu hrany.}\label{fig:decisionTreeOrientGraphs}
\end{figure}

\begin{def-sk}[Symetrizácia]\label{kap1:2.3:2.3.1:symetrizacia}
Každému orientovanému grafu $G = (V,E)$ môžeme priradiť neorientovaný graf $sym(G) = (V,E')$, kde $E' = \{\{x,y\}; (x,y) \in E$ alebo $(y,x) \in E\}$. Graf $sym(G)$ nazývame symetrizáciou grafu G.
\end{def-sk}

\begin{def-sk}[Súvislosť]\label{kap1:2.3:2.3.1:suvislost}
Hovoríme, že graf G je súvislý, keď pre každé dva jeho vrcholy x a y v ňom existuje cesta z x do y. Orientovaný graf je súvislý, keď je súvislá jeho symetrizácia (takáto súvislosť sa nazýva slabá).
\end{def-sk}

Pre orientované hrany existuje ešte druhý typ súvislosti a to silná. Avšak táto súvislosť je pre naše účely a hlavne definíciu rozhodovacích stromov nepodstatná.

\begin{def-sk}[Kružnica]\label{kap1:2.3:2.3.1:kruznica}
Graf $C_{n}$, kde n > 3, nazývame kružnicou, keď $ V = \{1,2,...,n\}$ a $E = \{\{i,i+1\};i=1,...,n-1\} \cup \{\{1,n\}\}$.
\end{def-sk}

Na obrázku \ref{fig:decisionTreeCycles} sú príklady štyroch druhov kružníc (tých najmenších).

\begin{figure}[h]
\centering
\centerline{\mbox{\includegraphics[width=300pt]{../img/kap1/DT-cycles.pdf}}}
\caption{Pár príkladov kružníc. Čierne body predstavujú vrcholy a čiary medzi vrcholmi reprezentujú hrany z $E$.}\label{fig:decisionTreeCycles}
\end{figure}

\begin{def-sk}[Strom]\label{kap1:2.3:2.3.1:strom}
Strom je súvislý graf neobsahujúci kružnicu. Orientovaný strom je súvislý orientovaný graf, ktorý po symetrizácii neobsahuje žiadnu kružnicu.
\end{def-sk}

Podrobnejší výklad o stromoch môže čitateľ nájsť v knihe \cite{kap1-KapitolyDiskretka}. 

\begin{def-sk}[Zakorenený strom]\label{kap1:2.3:2.3.1:korenovystrom}
Zakorenený strom je orientovaným stromom, v ktorom je jeden význačný vrchol nazývaný koreň stromu. Hrany takéhoto stromu sú jednoznačne orientované a vedú smerom od koreňa. Pre zakorenený strom ďalej definujeme tieto pojmy:
\begin{itemize}
\item \textbf{potomok} vrcholu $V$ je každý vrchol, do ktorého je vedená orientovaná hrana z vrcholu $V$,
\item \textbf{list} je vrchol, ktorý nemá ďalšieho potomka.
\end{itemize} 
\end{def-sk}

\begin{figure}[h]
\centering
\centerline{\mbox{\includegraphics[width=200pt]{../img/kap1/DT-parentchild.pdf}}}
\caption{Popis jednotlivých druhov vrcholov vzhľadom k danému vrcholu (červený). Modrá -- nasledovníci, žltá -- predchodcovia, zelená -- potomkovia, fialová -- rodič)}\label{fig:decisionTreeParentChild}
\end{figure}

\begin{def-sk}[Vzťahy v strome]\label{kap1:2.3:2.3.1:naslednik}
Nech $G = (V,E)$ je zakorenený strom (obrázok \ref{fig:decisionTreeParentChild}) a $k$ je koreň stromu, potom uvažujme:
\begin{itemize}
\item \textbf{predchodca} vrcholu $v$ je každý vrchol na orientovanej ceste z koreňa $k$ do vrcholu $v$,
\item \textbf{nasledovník} vrcholu $v$ je každý vrchol, kde jedným z predchodcov tohoto vrcholu je vrchol $v$,
\item \textbf{rodič} vrcholu $v$ je vrchol, z ktorého vedie orientovaná hrana do vrcholu $v$. Inak aj bezprostredný predchodca vrcholu $v$,
\item \textbf{potomok} vrcholu $v$ je vrchol, do ktorého vedie orientovaná hrana z vrcholu $v$. Inak aj bezprostredný nasledovník vrcholu $v$.
\end{itemize} 
\end{def-sk}

\subsection{Rozhodovací strom}\label{kap1:2.3:2.3.2}
Rozhodovací strom je prediktívny model, ktorý môžeme použiť pri riešení klasifikačných, ale aj regresných úloh. Z časti \ref{kap1:2.2:2.2.4:Supervised} vieme, že táto technika patrí pod metódy učenia s učiteľom. Z tejto časti taktiež vieme, že trénovacie dáta majú mierne pozmenený zápis od toho zadefinovaného v oddieli \ref{kap1:2.1:2.1.1:DataRepresentation}.

\begin{def-sk}[Trénovacie dáta]\label{kap1:2.3:2.3.2:traindata}
Trénovacie dáta sú n prvkovou množinou dvojíc $D = \{(\mathbf{x}_{k},y_{k})\}_{k=1}^{n}$, kde $\mathbf{x}_{k}$ je $d$-rozmerný príznakový vektor $d \geq 1$ a $y_{k}$ je jeho klasifikácia/regresia vybraná z množiny $C$. Pre klasifikáciu je doménou $C$ konečná množina symbolov, $C = \{c_{1},c_{2},\ldots,c_{m}\}$. Pri regresii je doména $C$ nejaká nekonečná množina, ako napríklad interval reálnych čísel $\mathbb{R}$, podmnožina celých čísel, $\ldots$.
\end{def-sk}

\begin{def-sk}[Podmnožina dát podľa hodnoty atribútu]\label{kap1:2.3:2.3.2:subsetValueData}
Nech $D = \{(\mathbf{x}_{i},y_{i})\}_{i=1}^{n}$ sú trénovacie dáta. Ďalej nech $D_{X_{i} = t}$ je podmnožina $D$ taká, že obsahuje práve všetky dvojice $(\mathbf{x},y) \in D$, pre ktoré platí $\mathbf{x} = (x_{1}, \ldots, x_{d})$ a $x_{i} = t$. Podobne $D_{c_{i}}$ je podmnožina dátovej množiny $D$ skladajúca sa z dvojíc $(\mathbf{x},y) \in D$, pre ktoré platí $y = c_{i}$.
\end{def-sk}

Na obrázku \ref{fig:decisionPlaneDataa} sú zobrazené štyri množiny dát, ktoré dokopy tvoria celú množinu $D$ ($=D_{1} \cup D_{2} \cup D_{3} \cup D_{4}$). Farba jednotlivých útvarov určuje ich klasifikáciu.
 
\begin{figure}[h]
\centering
\begin{subfigure}[b]{0.45\textwidth}
\includegraphics[width=\textwidth]{../img/kap1/DT-planedata.pdf}
\caption{}\label{fig:decisionPlaneDataa}
\end{subfigure}
\qquad
\begin{subfigure}[b]{0.45\textwidth}
\includegraphics[width=\textwidth]{../img/kap1/DT-planedatasep.pdf}
\caption{}\label{fig:decisionPlaneDatab}
\end{subfigure}
\caption{Umiestnenie dátových množín $D = D_{1} \cup D_{2} \cup D_{3} \cup D_{4}$ v rovine. Farby určujú klasifikáciu/regresiu pre množinu. V prípade b) sú vyznačené aj deliace nadroviny.}\label{fig:decisionPlaneData}
\end{figure}

\begin{def-sk}[Zobrazenie na atribúty]\label{kap1:2.3:2.3.2:mapovanievrcholov}
Nech P je množina príznakových vektorov v daných dátach, $G = (V,E)$ je zakorenený strom a L je množina všetkých listov tohoto stromu. Potom zobrazením $\alpha:(V \setminus L) \times P \rightarrow V$ priraďujeme nelistovému vrcholu $u \in V \setminus L$ a špecifickému príznakovému vektoru $\mathbf{x} \in P$ vrchol $v \in V$ (listový alebo nelistový), taký že $(u,v)$ je orientovaná hrana z $E$.
\end{def-sk}

\begin{def-sk}[Zobrazenie na výstupný atribút]\label{kap1:2.3:2.3.2:mapovanielistov}
Nech C je výstupný atribút a dom(C) jeho obor hodnôt v daných dátach, $G$ je koreňový strom a L je množina všetkých listov tohoto stromu. Potom zobrazenie $\beta:L \rightarrow dom(C)$ priraďuje každému listu triedu alebo reálnu hodnotu.
\end{def-sk}

\begin{def-sk}[Rozhodovací strom]\label{kap1:2.3:2.3.2:DT}
Majme zakorenený strom $G = (V,E)$ a zobrazenia $\alpha$ a $\beta$, ako sú uvedené v Definíciach \ref{kap1:2.3:2.3.2:mapovanielistov} a \ref{kap1:2.3:2.3.2:mapovanievrcholov}. Rozhodovací strom pre trénovacie dáta $D$ definujeme ako trojicu $(G,\alpha,\beta)$.
\end{def-sk}

Pri takejto definícii klasifikácia/regresia funguje nasledovne. Daný príznakový vektor je testovaný na nelistovom vrchole pomocou funkcie $\alpha$, ktorou získame ďalší vrchol na otestovanie. Tento proces začína v koreni stromu a pokračuje až dokým nedostaneme nejaký listový vrchol. V tomto vrchole, za použitia funkcie $\beta$ dostávame požadovanú klasifikáciu/regresiu.
Testovanie väčšinou prebieha porovnávaním určitého atribútu alebo skupiny atribútov vektora s nejakou hodnotou.

\begin{def-sk}[Deliaca nadrovina]\label{kap1:2.3:2.3.2:axisHyperplanes}
Deliacu nadrovinu $h$, určenú vektorom $\mathbf{w}$ a číslom $b$ definujeme ako množinu všetkých bodov $\mathbf{x}$, pre ktoré platí $h:\mathbf{w^{T}x} + b = 0$, kde $\mathbf{w^{T}}$ nazývame normálový váhový vektor nadroviny $h$ a $b$ je vzdialenosť nadroviny $h$ od počiatku súradnicovej sústavy.
\end{def-sk}

V modeli rozhodovacích stromov sú uvažované oddeľujúce nadroviny, ktoré sú kolmé na jednu z hlavných osí. Z Definície \ref{kap1:2.3:2.3.2:axisHyperplanes} teda vyplýva, že normálový váhový vektor $\mathbf{w}$ takejto nadroviny musí byť taktiež rovnobežný s niektorou z hlavných osí. Preto je možné tento vektor obmedziť iba na jeden z množiny $\{\mathbf{e}_1,\mathbf{e}_2,\ldots,\mathbf{e}_d\}$, kde každý vektor $\mathbf{e}_i \in \mathbb{R}^{d}$ a obsahuje, okrem jedničky na $i$-tom mieste, samé nuly. Potom pre každý bod $\mathbf{x} = \{x_{1},x_{2},\ldots,x_{d}\}$ nadroviny $h$ platí:
\begin{align}
h:\mathbf{e}_i\mathbf{x} + b &= 0 \nonumber\\
h:x_{i} + b &= 0 \label{kap1:2.3:2.3.2:eq1}
\end{align}

\begin{def-sk}[Deliaci bod]\label{kap1:2.3:2.3.2:splitPoints}
Deliaci bod zodpovedajúci rozhodnutiu a definovaný nadrovinou rozdeľuje dátový priestor $\mathcal{R}$ na dva podpriestory $\mathcal{R}_{p}$ a $\mathcal{R}_{n}$. 
\end{def-sk}

\begin{remark-sk}\label{kap1:2.3:2.3.2:remarkSplitPoints}
Takže všetky body $\mathbf{x}$, pre ktoré platí $h(\mathbf{x}) < 0$, sú na jednej strane nadroviny, pričom $h(\mathbf{x}) \geq 0$ implikuje opačnú stranu. V prípade $h(\mathbf{x}) < 0$ podľa (\ref{kap1:2.3:2.3.2:eq1}) platí, že $x_{i} + b < 0$ a teda $x_{i} < -b$, pričom $x_{i}$ nadobúda nejakú hodnotu z domény atribútu $X_{i}$. Deliaci bod v rozhodovacom strome je preto potom uvádzaný vo formáte $X_{i} < v$, kde $v = -b$.
\end{remark-sk}

Rozhodovacie stromy teda obsahujú vrcholy, ktoré reprezentujú rozhodnutia zodpovedajúce deliacim bodom asociovanými s deliacimi nadrovinami.

\begin{def-sk}[Binárny rozhodovací strom]\label{kap1:2.3:2.3.2:binarnyDT}
Binárny rozhodovací strom je rozhodovací strom, pre ktorý platí, že počet potomkov každého vrcholu tohoto stromu je rovný 2 alebo 0 (len v prípade listov).
\end{def-sk}

Na obrázku \ref{fig:decisionPlaneDatab} sú nakreslené 4 množiny bodov zo štyroch tried spolu s deliacimi nadrovinami. Rozhodovací strom, v konečnom dôsledku, svojou vnútornou stavbou pokladá tieto deliace čiary.

\section{Typy rozhodovacích stromov}\label{kap1:2.4:DTTypes}
Z časti taxonómia metód \ref{kap1:2.2:2.2.3:Taxonomy} vieme, že rozhodovací strom je predikčný model.
Uvádzali sme, že predikčné modely sú dvoch typov, klasifikačné a regresné. Delenie rozhodovacích stromov bude v tomto ohľade také isté. Informácie použité v tejto časti sme čerpali už zo spomínaných kníh \cite{kap1-DataMiningForTrees,kap1-DataMiningForMasses}, článku \cite{kap1-DecisionTreesTypes} a z online zdroja \cite{online-DecisionTreeTypes}.

Čo je to klasifikačný strom spolu s jeho dôležitými vlastnosťami zhrnieme v časti \ref{kap1:2.4:2.4.1:DTClassification}. Poznatky o regresných stromoch budú uvedené v časti \ref{kap1:2.4:2.4.1:DTRegression}. V obidvoch častiach ukážeme prevod rozhodovacieho stromu na súbor jednoduchých pravidiel pre ľahšiu predikciu.
\subsection{Klasifikačný strom}\label{kap1:2.4:2.4.1:DTClassification}
Rozhodovací strom, ako už vieme z predchádzajúcej kapitoly je hierarchická štruktúra, ktorá pomocou rozhodnutí dovedie zadaný vektor k určitej výstupnej hodnote. V obore dobývania znalostí je rozhodovací strom prediktívny model, ktorý môže reprezentovať ako klasifikačný, tak aj regresný model. Keď je strom použitý na klasifikačné úlohy, tak ho nazývame klasifikačný.

Klasifikačný strom používame vtedy, keď chceme klasifikovať objekt do triedy vybranej z konečnej, predom definovanej množiny tried podľa jeho príznakov. Klasifikačné stromy sú použiteľné najmä ako objaviteľská technika. Široko sú využívané v oblasti financií (pre ich zrozumiteľnosť), medicíny, vo vedeckých okruhoch, ale aj priemyselných procesoch. Aj napriek svojim výhodám nenahradzujú iné tradičné techniky predikcie ako sú Bayesovské alebo neurónové siete. Pri klasifikačných stromoch si je taktiež nutné uvedomiť, že atribúty vo vrcholoch môžu byť ako numerické, tak aj kategorické.

Na Obrázku \ref{fig:TypesClassify} je nakreslený typický klasifikačný strom. Tento príklad vznikol priamym prepisom obrázku \ref{fig:decisionPlaneDatab} spolu s uvedenými deliacimi nadrovinami. Uvažujme, že každej definovanej množine $D_{1},\ldots,D_{4}$ z \ref{fig:decisionPlaneDatab} priradíme triedu $c_{1},\ldots,c_{4}$. Podľa pozorovania \ref{kap1:2.3:2.3.2:remarkSplitPoints} z predchádzajúcej časti je vidieť, že testy atribútov vo vnútorných uzloch stromu zodpovedajú deliacim bodom a konečné klasifikácie jednotlivým dátovým množinám na tomto obrázku. Vstupnú inštanciu, ktorú chceme klasifikovať označme $I$. Pravidlá tohoto klasifikačného stromu vyzerajú takto:
\begin{align}
Y(I) < 3.5 \wedge X(I) < 3.5 &\Rightarrow c_{2} \nonumber \\
Y(I) < 3.5 \wedge X(I) \geq 3.5 &\Rightarrow c_{3} \nonumber \\
Y(I) \geq 3.5 \wedge X(I) > 2 &\Rightarrow c_{4} \nonumber \\
Y(I) \geq 3.5 \wedge X(I) \leq 2 &\Rightarrow c_{1} \nonumber
\end{align}
\begin{figure}[h]
\centering
\centerline{\mbox{\includegraphics{../img/kap1/DT-binary.pdf}}}
\caption{Príklad klasifikačného stromu (binárneho), ktorý popisuje deliace nadroviny z obrázku \ref{fig:decisionPlaneDatab}.}\label{fig:TypesClassify}
\end{figure}
Na tvorbu klasifikačných stromov sa využíva mnoho rôznych techník. 
Prvý publikovaný algoritmus sa nazýva THeta Automatic Interaction Detection (THAID). Neskôr vznikli vylepšené techniky ako Iterative Dichotomiser 3 (ID3), C4.5 (zlepšený ID3), CHi-squared Automatic Interaction Detection (CHAID) a Classification And Regression Tree (CART), ktorý sa dá použiť pri tvorbe klasifikačných aj regresných stromov. Každý z nich používa iný typ hľadania deliacich bodov. Medzi ďalšie, menej známe techniky môžeme uviesť CRUISE, GUIDE, QUEST, o ktorých si je možné prečítať v článkoch \cite{kap1-DecisionTreesUnused1,kap1-DecisionTreesUnused2,kap1-DecisionTreesUnused3}.
Konkrétnejšie budú niektoré z nich popísané v oddieli \ref{kap1:2.5:DTTechniques}. 

\subsection{Regresný strom}\label{kap1:2.4:2.4.1:DTRegression}
Regresný strom je variantou rozhodovacieho stromu navrhnutý na aproximáciu reálnych funkcií na rozdiel od klasifikácie. Strom je veľmi podobný tomu klasifikačnému. Jediným rozdielom je výstupný atribút, ktorého doména je v tomto prípade nekonečná množina ($\mathbb{R},\mathbb{N},\ldots$).

\begin{figure}[h]
\centering
\centerline{\mbox{\includegraphics{../img/kap1/DT-regression.pdf}}}
\caption{Príklad regresneho stromu (binárneho), ktorý predikuje požičky pre žiadateľov.}\label{fig:TypesRegression}
\end{figure}

Tento druh stromov teda používame vtedy, keď chceme predikovať vlastnosť objektu, ktorá je z nekonečnej množiny. Využitie je úplne rovnaké, ako pri klasifikačných stromoch. Regresné stromy majú ale bohatšie spôsoby pri predikcii. Týka sa to hlavne ďalšieho využitia iných druhov modelov v koncových listoch. Takto vytvorené stromy bývajú ale o niečo náročnejšie na konštrukciu kvôli ich komplikovanejšej štruktúre. Vďaka zložitejšej štruktúre stromy dávajú pri regresii zvyčajne lepšie výsledky než iné jednoduchšie metódy.

Na obrázku \ref{fig:TypesRegression} si je možné všimnúť typický regresný strom. V tomto prípade predikuje výšku priradenej pôžičky žiadateľom podľa ich veku a sociálneho zaradenia (príklad nevychádza z reálnych údajov). Označme žiadateľa/entitu o pôžičku ako $E$. Jednotlivé pravidlá pre tento strom by vyzerali takto:
\begin{align}
vek(E) < 25 \wedge soc.zar(E) = \check{s}tudent &\Rightarrow schv\acute{a}len\acute{a}\  po\check{z}i\check{c}ka = 30000, \nonumber \\
vek(E) < 25 \wedge soc.zar(E) \neq \check{s}tudent &\Rightarrow schv\acute{a}len\acute{a}\  po\check{z}i\check{c}ka = 15000, \nonumber \\
vek(E) \geq 25 \wedge soc.zar(E) = zamestnan\acute{y} &\Rightarrow schv\acute{a}len\acute{a}\  po\check{z}i\check{c}ka = 60000, \nonumber \\
vek(E) \geq 25 \wedge soc.zar(E) = nezamestnan\acute{y} &\Rightarrow schv\acute{a}len\acute{a}\ po\check{z}i\check{c}ka = 5000. \nonumber 
\end{align}
Regresné stromy sa konštruujú rôznymi technikami. Historicky prvým algoritmom na konštrukciu regresných stromov bol algoritmus Automatic Interaction Detection (AID), ktorý sa objavil pár rokov pred THAID. Ďalším, veľmi obľúbeným algoritmom je už spomínaný CART, ktorý dokáže vytvoriť popri regresných stromov aj tie klasifikačné. K menej známym technikám patria M5' a GUIDE.

\section{Techniky tvorby stromov}\label{kap1:2.5:DTTechniques}
Cieľom tejto časti je predviesť najznámejšie techniky tvorby stromov (regresných aj klasifikačných). Tieto techniky spadajú pod oblasť učenia s učiteľom, ktorá bola bližšie popísaná v oddieli \ref{kap1:2.2:2.2.4:Supervised}. Jedná sa o algoritmy, ktoré vytvárajú strom z trénovacej množiny postupne od koreňa. Dáta sú v algoritme rozdelené podľa deliaceho kritéria na podmnožiny. Ďalej pokračujeme rekurzívnym volaním toho istého algoritmu na každú z týchto podmnožín. V obore dobývania znalostí sa celý tento proces nazýva indukcia rozhodovacích stromov zhora nadol, anglicky top down induction of decision trees (TDIDT). Indukčné algoritmy sú príkladom hladných techník. V praxi sú jedným z najpoužívanejších stratégii pri trénovaní rozhodovacích stromov. V tejto časti sme čerpali z kníh \cite{kap1-DataMiningForTrees,kap1-DataMiningAndAnalysis} pričom na rozšírenie ďalších poznatkov sme použili online zdroje \cite{online-SplitCriterias,online-SplitCriteriasMatter,online-DTLectures}.

Najskôr v časti \ref{kap1:2.5:2.5.1:SplitCriterias} zadefinujeme, čo sú to kritéria delenia a vypíšeme tie najznámejšie z nich. V poslednej časti \ref{kap1:2.5:2.5.4:Generic} ukážeme generický algoritmus na indukciu stromov. 
\subsection{Kritéria delenia}\label{kap1:2.5:2.5.1:SplitCriterias}
Kritéria delenia sú funkcie používané pri konštrukcii stromov. Tieto kritéria majú za úlohu zvoliť vhodný atribút do vnútorného vrcholu stromu.
V rozhodovacích stromoch majú tieto vnútorné vrcholy väčšinou jediný atribút, na ktorom porovnávame jednotlivé inštancie z dát. Kritéria, ktoré vytvárajú takýto typ atribútov nazývame jednorozmerné.

Jednorozmerné kritéria môžu byť ďalej charakterizované podľa
\begin{itemize}
\item pôvodu (teória informácii, závislosť, vzdialenosť),
\item štruktúry, kam patria binárne kritéria, kritéria založené na miere neusporiadanosti (impurity-based) a ich normalizované verzie.
\end{itemize}

Viac príkladov kritérií, doplnené informácie o spomínaných kritériách a podrobnejšie vysvetlenie deliacich kritérií je možné nájsť v knihách \cite{kap1-DataMiningForTrees,kap1-StatisticLearn}. 
\subsubsection{Neusporiadanosť}\label{kap1:2.5:2.5.1:Impurity}
Na bližšie zadefinovanie kritérií založených na neusporiadanosti potrebujeme vedieť pár pojmov z teórie informácii.  
\begin{def-sk}[Miera neusporiadanosti]\label{kap1:2.5:2.5.1:ImpurityFunction}
Nech $k \in \mathbb{N}$. Miera/funkcia neusporiadanosti $\gamma: [0,1]^{k}  \rightarrow \mathbb{R}$ je funkcia definovaná pre $k$-tice $P = (p_{1},\ldots,p_{k})$, spĺňajúca $\sum_{i=1}^{k}p_{i} = 1, p_{i} \geq 0$ taká, že:
\begin{itemize}
\item $\gamma(P) \geq 0$,
\item $\gamma(P)$ dosahuje minimum v bodoch, pre ktoré platí, že $p_{i} = 1$, pre nejaké $i$, $1 \leq i \leq k$,
\item $\gamma(P)$ dosahuje maxima v bode $(1/k,\ldots,1/k)$,
\item $\gamma(P)$ je symetrická vzhľadom k hodnotám $p_{1},\ldots,p_{k}$,
\item $\gamma(P)$ je všade diferencovateľná.
\end{itemize}
\end{def-sk}

\begin{def-sk}[Neusporiadanosť množiny]\label{kap1:2.5:2.5.1:ImpuritySet}
Nech $D$ je dátová množina a $C = (c_{1},\ldots,c_{k})$ je množina klasifikačných tried. Ďalej nech $D_{c_{i}}$ je podmnožina dátovej množiny $D$ tvorená vektormi z triedy $c_{i}$, pre $i$=$1,\ldots,k$ a $\gamma: [0,1]^{k}  \rightarrow \mathbb{R}$ je funkcia neusporiadanosti. Potom neusporiadanosť dátovej množiny $D$, označovaná $\gamma(D)$, je definovaná ako
\begin{equation}
\gamma(D) = \gamma
	\left(
	\dfrac
		{\lvert D_{c_{1}}\lvert}
		{\lvert D \lvert},		
	\ldots,
	\dfrac
		{\lvert D_{c_{k}}\lvert}
		{\lvert D \lvert}				
	\right).
\end{equation}
\end{def-sk}

\begin{def-sk}[Zníženie neusporiadanosti množiny]
Nech $X_{i}$ je daný atribút a $D_{1}, \ldots, D_{n}$ sú podmnožiny dátovej množiny $D$, ktoré vznikli rozdelením $D$ podľa atribútu $X_{i}$. Potom výsledné zníženie neusporiadanosti označené ako $\Delta\gamma(X_{i},D) = \gamma(D) - \sum_{i=1}^{s}\dfrac{\lvert D_{i} \lvert}{\lvert D \lvert} \gamma(D_{i})$.
\end{def-sk}

\begin{figure}[h]
\centering
\centerline{\mbox{\includegraphics[width=400pt]{../img/kap1/DT-splitcrit.pdf}}}
\caption{Príklad, ako môžu kritéria delenia rozdeliť dátovú množinu so štyrmi triedami. Farby určujú pomery dát s danou triedou. Obrázok a) oproti b) je vhodnejší, pretože minimalizuje neusporiadanosť v listových vrcholoch týchto stromov.}\label{fig:SplitCriterias}
\end{figure}

Pri indukcii rozhodovacích stromov zvyčajne využívame nejakú hladnú stratégiu, ktorá sa snaží čo najviac znížiť neusporiadanosť celého rozhodovacieho stromu. Na obrázku \ref{fig:SplitCriterias} môžeme vidieť rozdelenie stromu v jednom z jeho vrcholov. V tomto obrázku preferujeme prípad a), pretože minimalizuje neusporiadanosť podmnožín $D$ v listových vrcholoch.
\subsubsection{Informačný zisk}\label{kap1:2.5:2.5.1:InfoGain} 
Informačný zisk je jedno z kritérií delenia založené na neusporiadanosti. Definícia kritéria vychádza z predchádzajúcej časti \ref{kap1:2.5:2.5.1:Impurity}, kde za funkciu neusporiadanosti $\gamma$ je dosadená entropia (z teórie informácii). Vzorec kritéria je potom definovaný ako ($D_{X_{i}}$ je v Definícii \ref{kap1:2.3:2.3.2:subsetValueData})

\begin{align}
\infogain(& X_{i},D) = \nonumber \\
& \entropy(C,D) -
\sum_{t \in dom(X_{i})}^{}
\dfrac{\lvert D_{X_{i} = t}\lvert}{\lvert D \lvert} 
\entropy(C,D_{X_{i} = t}). \label{kap1:2.5:2.5.1:InfoGainDef} 
\end{align}
kde
\begin{equation}
\entropy(C,D) = \sum_{c_{i} \in C}^{} \left( -
\dfrac{\lvert D_{c_{i}}\lvert}{\lvert D \lvert} \log
\dfrac{\lvert D_{c_{i}}\lvert}{\lvert D \lvert}\right). 
\end{equation}

Informačný zisk blízko súvisí s metódou maximálnej vierohodnosti. Táto metóda je populárnou hlavne v oblasti štatistiky pri odhadovaní parametrov nejakého štatistického modelu.

Ďalšie informácie o entropii je možné nájsť v knihách spomenutých v týchto častiach. Pre naše potreby nám stačí použiť Wikipédiu \cite{wiki-Entropy}.
\subsubsection{Gini Index}\label{kap1:2.5:2.5.1:GiniIndex}
Gini index je ďalšie z kritérii založené na neusporiadanosti. Používa sa ako miera diverzifikácie hodnôt vo výstupnom atribúte. Tak ako aj ostatné miery neusporiadanosti nadobúda hodnoty od 0 po 1, kde 0 znamená dokonalú rovnosť tried výstupu a 1 úplnú nerovnosť. Kritérium využívajúce Gini index ako funkciu neusporiadanosti $\gamma$ potom definujeme ako
\begin{align}
\ginindex(X_{i},D) = \gini(C,D) -
\sum_{t \in dom(X_{i})}^{}
\dfrac{\lvert D_{X_{i} = t}\lvert}{\lvert D \lvert} 
\gini(C,D_{X_{i} = t}). \label{kap1:2.5:2.5.1:GiniIndexDef}
\end{align}
kde
\begin{equation}
\gini(C,D) = 1 - \sum_{c_{i} \in C}^{} 
\left(
\dfrac{\lvert D_{c{i}}\lvert}{\lvert D \lvert} 
\right) ^ 2. 
\end{equation}
Taktiež je možné nahliadnuť, že pre binárny prípad (počet tried je rovný dvom) je výpočet kritéria jednoduchší.
\begin{equation}
\gini(C,D) = 2p_{c_{1}}p_{c_{2}}, 
\end{equation}

kde $p_{c_{i}}$ je relatívny počet inštancií, ktorých výstupný atribút má hodnotu $c_{i}$.

\subsubsection{Koeficient zisku}\label{kap1:2.5:2.5.1:GainRatio} 
Koeficient zisku je ďalším známym kritériom, ktoré pre svoj výpočet využíva informačný zisk. Patrí medzi normalizované kritéria neusporiadanosti. Definované je ako
\begin{align}
\gainratio(X_{i},D) =
\dfrac{\infogain(X_{i},D)}{\entropy(X_{i},D)} \label{kap1:2.5:2.5.1:GainRatioDef} 
\end{align}

Z definície je vidieť, že kritérium nie je definované, keď je entropia rovná nule. Taktiež je kritérium náchylné k atribútom, pre ktoré je entropia veľmi malá. V minulosti bolo ukázané, že toto kritérium dokázalo prekonať normálny informačný zisk v presnosti aj v zložitosti.

\subsubsection{Twoing kritérium}\label{kap1:2.5:2.5.1:Twoing} 
Medzi menej známe kritéria patrí Twoing kritérium \cite[s.88]{kap1-DataMiningForTrees}. Je to binárne kritérium, ktoré je ako zbytok binárnych kritérii založené na rozdelení domény nejakého atribútu na dve podmnožiny. Zavedené bolo hlavne kvôli Gini indexu a jeho problému s príliš veľkou doménou výstupného atribútu. Kritérium je definované nasledovne 
\begin{align}
\twoing(X_{i},& dom_{1}(X_{i}),dom_{2}(X_{i}),D) = \nonumber \\
& 0.25 
\dfrac{\lvert D_{X_{i} = t \in dom_{1}(X_{i})}\lvert}{\lvert D\lvert}
\dfrac{\lvert D_{X_{i} = t \in dom_{2}(X_{i})}\lvert}{\lvert D\lvert} \nonumber \\
& \left(
\sum_{c_{i} \in C}^{}
\left| 
\dfrac{
\lvert D_{X_{i} = t \in dom_{1}(X_{i})} \cap D_{c_{i}}\lvert }{\lvert D_{X_{i} = t \in dom_{1}(X_{i})}\lvert} -
\dfrac{
\lvert D_{X_{i} = t \in dom_{2}(X_{i})} \cap D_{c_{i}}\lvert }{\lvert D_{X_{i} = t \in dom_{2}(X_{i})}\lvert}
\right|
\right)^2, \label{kap1:2.5:2.5.1:TwoingDef} 
\end{align}

kde $dom_{1}(X_{i})$ a $dom_{2}(X_{i})$ sú podmnožiny domény atribútu $X_{i}$. Keď je výstupný atribút binárny, tak Gini Index a Twoing kritérium sú totožné. Toto kritérium v nebinárnom prípade preferuje atribúty s podobným rozdelením dátovej množiny. Algoritmus CART používa toto deliace kritérium pri indukcii stromov.

\subsection{Vyhodnotenie kvality stromov}\label{kap1:2.5:2.5.2:Evaluation}

\subsection{Orezávanie stromov}\label{kap1:2.5:2.5.3:Pruning}

\subsection{Generický indukčný algoritmus}\label{kap1:2.5:2.5.4:Generic}
Indukčné algoritmy majú za úlohu vytvoriť rozhodovací strom zo zadanej trénovacej množiny $D$. Hlavnou úlohou pri tvorbe je minimalizovať generalizačnú chybu. Nie vždy sa jedná len o minimalizovanie. Algoritmus sa môže snažiť o vytvorenie stromu s čo najnižšou výškou alebo tiež s obmedzeným počtom vrcholov.

Tvorba optimálneho rozhodovacieho stromu je ťažkou úlohou. Viac o tom v knihe \cite[s.51]{kap1-DataMiningForTrees}. Indukčné algoritmy sú dvoch typov. Tie, ktoré vytvárajú stromy zhora nadol a zdola nahor. Medzi algoritmy typu zhora nadol patria ID3, C4.5, CART a mnoho ďalších. Tvorba stromu je pri niektorých algoritmoch (C4.5 a CART) spojená aj s jeho orezaním. Väčšinou na začiatku existuje jeden vrchol stromu (koreň), ktorému odpovedá celá trénovacia množina $D$. Následníci tohoto koreňa ďalej zodpovedajú podmnožinám trénovacej množiny. Týmto postupom stále delíme trénovaciu množinu, a tým vytvárame nové vrcholy stromu. Delenie prebieha až dokým nie je pravdivé niektoré zo stanovených kritérií zastavenia. Kostra takejto tvorby stromu je popísaná v Algoritme  \ref{fig:genericAlgoritm}.

\begin{algorithm} 
\floatname{algorithm}{Algoritmus}
\caption{Generický algoritmus na tvorbu stromov, z ktorého vychádzajú známe algoritmy ID3,C4.5,a pod.}\label{fig:genericAlgoritm}
$D$ - Trénovacia množina \\
$X$ - Množina atribútov \\
$Y$ - Výstupný atribút \\
$Del\_Krit$ - Kritérium pre delenie \\
$Stop\_Krit$ - Kritérium na zastavenie tvorby stromu \\
$T$ - Vytvorený strom 
\begin{algorithmic}
\Function{IndukciaStromu}{$D$,$X$,$Y$,$Del\_Krit$, $Stop\_Krit$}
\State $\var{T} \gets $ Strom s jedným vrcholom (koreň)
\If {$Stop\_Krit(D)$}:	  
\State $\var{T} \gets $ list s najčastejšou hodnotou atribútu $Y$ v $D$.
\Else
\State $\var{A} \gets min_{X_{i}}Del\_Krit(X_{i},D)$
\State $attr(\var{T}) \gets A$
\EndIf
\ForAll{$\var{v_{i}} \in dom(A)$}:
\State $PodStrom_{i}(T) \gets $ \parbox[t]{280pt}{\textsc{IndukciaStromu}($D_{X_{i} = v_{i}}$,$X$,$Y$,$Del\_Krit$, $Stop\_Krit$)} 
\State $Hrana_{i}(T) \gets v_{i}$  
\EndFor \\
\Return{\textsc{OrezanieStromu}($D$,$T$,$Y$)}
\EndFunction
\\
\Function{OrezanieStromu}{$D$,$T$,$Y$}
\Repeat 
\State $t \gets$ \parbox[t]{350pt}{textvrchol stromu T, ktorý po orezaní maximalizuje predom zvolené vyhodnocovacie kritérium}
\If{$t \neq \emptyset$}:
\State $T \gets$ orezanie($T$,$t$)
\EndIf
\Until{$t = \emptyset$} \\
\Return{$T$}
\EndFunction
\end{algorithmic}
\end{algorithm}

%\begin{figure}[h]
%\lstset{inputencoding=utf8,extendedchars=true,basicstyle=\ttfamily\small,literate={á}{{\'a}}1 {é}{{\'e}}1 {ž}{{\v{z}}}1 {ň}{{\v{n}}}1 {š}{{\v{s}}}1 {č}{{\v{c}}}1 {ď}{{\v{d}}}1 {í}{{\'i}}1 {ý}{{\'y}}1 {ú}{{\'u}}1}
%\begin{lstlisting}[mathescape]
%IndukciaStromu($D$,$A$,$y$,$DeliaceKrit\acute{e}rium$,$Krit\acute{e}riumZastavenia$)
%$D$ - Trénovacia množina
%$A$ - Vstupné príznaky
%$y$ - Výstupný atribút
%$DeliaceKrit\acute{e}rium$ - Kritérium pre delenie
%$Krit\acute{e}riumZastavenia$ - Kritérium na zastavenie tvorby stromu
%
%Vytvor strom $T$ s jedným vrcholom (koreň)
%IF Krit\acute{e}riumZastavenia(D) THEN
%	Sprav z $T$ list s najčastejšou
%	hodnotou atribútu $y$ v $D$.
%ELSE
%	Vyber atribút $a \in A$ s najlepšou
%	hodnotou $DeliaceKrit\acute{e}rium(a,D)$.
%	Aktuálnemu vrcholu priraď atribút $a$.
%	FOR $v_{i} \in A$:
%		Nastav $PodStrom_{i}$ = IndukciaStromu($D_{v_{i}}$,$A$,$y$,_,_).
%		Spoj vrchol $t_{T}$ a $PodStrom_{i}$ s hranou označenou
%		$v_{i}$.
%	END FOR
%END IF
%RETURN OrezanieStromu(D,T,y).
%\end{lstlisting}
%\caption{Generický algoritmus z ktorého vychádzajú známe algoritmy ID3,C4.5,$\ldots$. Prevzaté a prepísané z \cite[s.52]{kap1-DataMiningForTrees}}\label{fig:genericAlgoritm}
%\end{figure}

Existujú rôzne kritéria zastavenia, ktoré prerušia vytváranie stromu. Tými najznámejšími sú

\begin{itemize}
\item Všetky inštancie v trénovacej množine majú jednotný výstupný atribút $y$.
\item Strom dosiahol maximálnu výšku.
\item Hodnota najlepšieho deliaceho kritéria je menšia než určitý prah.
\item Počet inštancií v danom uzle je menší než dopredu zadaná konštanta.
\end{itemize}

V ďalších častiach rozoberieme tri najznámejšie indukčné algoritmy na tvorbu stromov. Konkrétne v časti \ref{kap1:2.5:2.5.4:ID3} ukážeme algoritmus ID3. Následne v časti \ref{kap1:2.5:2.5.4:C4.5} predstavíme vylepšený algoritmus C4.5. Pre regresné typy stromov predvedieme v časti \ref{kap1:2.5:2.5.4:CART} obľúbený algoritmus CART.

\subsubsection{ID3}\label{kap1:2.5:2.5.4:ID3}
ID3 algoritmus je najznámejší a najjednoduchší z indukčných algoritmov. Deliacim kritériom je Informačný zisk (\ref{kap1:2.5:2.5.1:InfoGainDef}). Kritériom zastavenia je nulový informačný zisk alebo dosiahnutie homogénnosti výstupného atribútu. Algoritmus nevyužíva žiadne techniky orezávania, nedokáže riešiť problémy s chýbajúcimi hodnotami v atribútoch a nevie pracovať s numerickými atribútmi.

Najväčšou výhodou ID3 je jeho jednoduchosť. Zvyčajne je kvôli tomu využívaný pre vzdelávacie účely. 

Nevýhody tohoto algoritmu presahujú spomenuté výhody. Stromy vytvorené týmto algoritmom sa nachádzajú v nejakom lokálnom optime (hladný algoritmus).
Vytvorené stromy bývajú väčšinou malé, ale keď algoritmus vytvorí väčší strom tak dochádza veľmi často k preučeniu. Preto pri tomto druhu algoritmu sú preferovanejšie menšie stromy. Algoritmus nepodporuje numerické hodnoty a pre prácu s nimi je nutné predspracovať dáta do nejakých diskrétnych skupín (hodnoty menšie ako $x$ budú v skupine $a$, ostatné v skupine $b$).
\subsubsection{C4.5}\label{kap1:2.5:2.5.4:C4.5}
C4.5 je vylepšením algoritmu ID3. Obidva majú rovnakého autora. Algoritmus používa ako deliace kritérium koeficient zisku (\ref{kap1:2.5:2.5.1:GainRatioDef}). Vytváranie stromu končí, keď počet inštancií pri delení je menší než zvolená konštanta. Zároveň využíva techniky orezávania a dokáže pracovať aj s numerickými atribútmi. C4.5 vie popri tomu  pracovať aj s chýbajúcimi hodnotami atribútov.

Indukcia pomocou C4.5 poskytuje oproti ID3 veľa zlepšení. Pri orezávaní odstraňuje vetvy, ktoré neprispievajú k presnosti a nahradzuje ich listami. Inštanciám môžu chýbať niektoré hodnoty atribútov. Numerické atribúty sú riešené rozdelením domény daného atribútu na dve podmnožiny podľa deliaceho bodu. Tento deliaci bod je zvolený tak, aby maximalizoval koeficient zisku.

Jedným z ďalších vylepšení C4.5 je ešte o niečo novší, komerčný algoritmus C5.0. Popri výhodám z C4.5 obsahuje ešte ďalšie zlepšenia, ako je rýchlosť výpočtu a zlepšenie práce s pamäťou. Ďalším zlepšením je vstavané využitie techniky boosting, ktorá dokáže zvýšiť prediktívne schopnosti modelu.

Implementovanú verziu algoritmu C4.5, nazývanú J48 môžeme nájsť už v spomínanom nástroji Weka.

\subsubsection{CART}\label{kap1:2.5:2.5.4:CART}
CART je hladný algoritmus, ktorý slúži na tvorbu klasifikačných aj regresných stromov. Autorom algoritmu je Breiman. Vytvorený strom je v tomto prípade vždy binárny (každý vnútorný vrchol ma dvoch potomkov). Algoritmus využíva Twoing kritérium (\ref{kap1:2.5:2.5.1:TwoingDef}) ako deliace kritérium. CART používa na orezávanie stromu metriku, ktorá berie do úvahy počet listov spolu s chybovosťou modelu (\cite{wiki-costcompprune}, \cite[s.382]{kap1-DecisionTree}).

Dôležitou vlastnosťou algoritmu je vytváranie regresných stromov. Vytvorené stromy obsahujú v listoch reálne hodnoty. V tomto prípade sú deliace body vyberané tak, aby minimalizovali kvadratickú odchýlku predikovaných hodnôt od správnych hodnôt na dátovej množine. V listových vrcholoch je hodnota rovná váženému priemeru výstupnej hodnoty dát v tomto vrchole.

\section{Využitie rozhodovacích stromov pri predikcii}\label{kap1:2.6:DTUsage}
V tomto oddieli zhrnieme výhody a nevýhody rozhodovacích stromov, ktoré priamo ovplyvňujú využiteľnosť tohoto modelu v praxi. Zároveň predstavíme projekty, ktoré prakticky využívajú techniku rozhodovacích stromov. Knihy s informáciami, ktoré boli použité v tejto časti sú \cite{kap1-DataMiningAndAnalysis,kap1-DataMiningForTrees}. Informácie s praktickým použitím rozhodovacích stromov sme čerpali zo stránok príslušných projektov \cite{online-astronomy} a \cite{online-psychoterapy}.

V časti \ref{kap1:2.6:2.6.1:AdvAndDis} popíšeme výhody a nevýhody rozhodovacích stromov, pri ich využití v praxi. V ďalších častiach predstavíme dva projekty, ktoré využívajú rozhodovacie stromy pri predikcii. Filtrovanie šumu v obrázkoch získaných z Hubblovho vesmírneho ďalekohľadu bude popísané v časti \ref{kap1:2.6:2.6.2:Hubble}. Využitie rozhodovacích stromov pre klinickú prax predvedieme v poslednej časti \ref{kap1:2.6:2.6.3:Clinical}.

\subsection{Výhody a nevýhody stromov}\label{kap1:2.6:2.6.1:AdvAndDis}
V predchádzajúcich oddieloch sme spomínali rôzne výhody a nevýhody rozhodovacích stromov. Medzi výhody môžeme zaradiť:
\begin{itemize}
\item Rozhodovacie stromy majú jednoduchú reprezentáciu, z ktorej je vidieť chovanie stromu. Pri malom počte listov je reprezentácia ľahko spracovateľná aj neprofesionálnym užívateľom.
\item Model je dostatočne bohatý na to, aby sa naňho dal previesť ľubovoľný iný diskrétny klasifikátor.
\item Cesty od koreňa k listom predstavujú jednotlivé pravidlá. Existuje teda jednoduchý prevod medzi stromom a množinou pravidiel.
\item Model rozhodovacích stromov je zaraďovaný medzi neparametrické techniky a teda neuvažuje žiadne predpoklady o štruktúre klasifikátora či rozdelení dát.
\item Rozhodovacie stromy dokážu pracovať s dátami, ktoré obsahujú chyby alebo ktorým chýbajú hodnoty niektorých atribútov.
\item Rozhodovacie stromy vedia pracovať s kategoriálnymi atribútmi, tak ako aj s numerickými.
\end{itemize}

K nevýhodám rozhodovacích stromov patrí:
\begin{itemize}
\item Hladné techniky tvorby stromov môžu uviaznuť v lokálnom optime. Stromy sú taktiež náchylné na preučenie.
\item Tvorba stromu hladnými technikami sa môže zamerať na nerelevantné atribúty a šum v dátach.
\item Krátkozrakosť algoritmu, kde delenia v jednotlivých vrcholoch pracujú iba s priamymi potomkami vrcholu.
\item Väčšina algoritmov požaduje, aby bol výstupný atribút kategoriálny (C4.5, ID3, \ldots).
\item Metódy, ktorými sa riešia problémy chýbajúcich hodnôt niektorých atribútov. Trénovacie dáta s chýbajúcimi hodnotami môžu byť napríklad odstránené z trénovacej množiny. Keď hodnota atribútu nechýba náhodou a nevyplnenie popisuje nejakú vlastnosť, tak je odstraňovanie inštancií úplne nevhodné (napr. neuvedený plat u dátovej množiny osôb môže znamenať, že osoba má príliš vysoký/nízky plat a nechcela ho uviesť).
\item Problém fragmentácie pri rozdeľovaní dát vo vrcholoch. K tomuto problému dochádza v prípade, keď pri deleniach dátových množín vo vrcholoch rozdelíme dátové množiny na rovnako veľké podmnožiny. Strom, ktorý by takto rozdeľoval dátovú množinu veľkosti $n$, obsahuje na ceste z koreňa do listu iba $\log n$ vrcholov, čo je ekvivalentné s otestovaním maximálne $\log n$ atribútov. Toto býva veľký problém pri existencii väčšieho počtu významných atribútov.
\end{itemize}

\subsection{Filtrovanie šumu z Hubblovho ďalekohľadu}\label{kap1:2.6:2.6.2:Hubble}
Cieľom autorov tohoto projektu \cite{online-astronomy} bolo vytvoriť klasifikačný model, ktorý by dokázal na obrázkoch získaných z Hubblovho vesmírneho ďalekohľadu rozlišovať medzi hviezdami a šumom. Šum, ktorý je pri takto získaných fotkách veľmi častý, je spôsobený kozmickým žiarením. 

Trénovacie dáta pozostávajú z príznakov $x_{1},\ldots,x_{20}$. Jednotlivé inštancie získali z okolia (veľkosti $3 \times 3$ pixele) zameraného na svetelné úkazy a z poznatkov o rôznych druhoch žiarenia. Výslednú trénovaciu množinu použili pri konštrukcii rozhodovacie stromu, ktorým bol predikovaný typ žiarenia -- dokázali rozlíšiť hviezdne žiarenie od šumu. Strom, ktorý vytvorili nazvali OC1. Model bol atypický v tom, že uvažoval aj take rozdeľovacie nadroviny, ktoré nie sú rovnobežné s hlavnými osami. Viac o algoritme je priamo v článku \cite[str. 281]{online-astronomy}.

\begin{figure}[h]
\centering
\centerline{\mbox{\includegraphics{../img/kap1/DT-hubble.pdf}}}
\caption{Získaný model klasifikujúci vstupný vektor s 20 príznakmi (Čísla v elipsách znamenajú delenia vo vrcholoch. Percentá pod vrcholmi predstavujú relatívny počet hviezd ku kozmickému žiareniu). Obrázok modelu je prebraný a prepísaný z \cite[str. 284]{online-astronomy}}\label{fig:HubbleTree}
\end{figure}

Na testovacích dátach dosahovali modely, vytvorené metódou OC1, presnosť až 95\%. Vytvorené stromy boli navyše kompaktné. Príklad takto vytvoreného stromu je na obrázku \ref{fig:HubbleTree}. Experimenty, nachádzajúce sa v článku uvádzajú, že zlepšením metód na eliminovanie šumu môže ďalej vylepšiť presnosť vytvorených klasifikátorov.

\subsection{Rozhodovacie stromy pre klinickú prax}\label{kap1:2.6:2.6.3:Clinical}
Cieľom práce bolo uľahčiť prácu doktorom a zdravotnému personálu vytvorením modelu, ktorý by varoval pred neúspechom liečby a pomocou správnych rozhodnutí odporúčal, ako čo najlepšie predísť takémuto neúspechu. 

Dáta, pre účely vytvorenia modelu, boli získané z iného projektu (Project TR-EAT), ktorý skúmal vzťahy medzi dĺžkou/intenzitou liečby poruchy prijímania potravy a výsledkom liečby. Z dát bola vybratá vzorka ľudí diagnostikovaná bulímiou nervosa ($n$ = 630). Pri výbere dát využívali hlavne údaje o prijatí, liečebnú stratégiu a odpoveď na liečbu. Na výslednú trénovaciu množinu bol použitý algoritmus CART.

\begin{figure}[h]
\centering
\centerline{\mbox{\includegraphics[width=385pt]{../img/kap1/DT-clinical.pdf}}}
\caption{Získaný model klasifikujúci výsledok liečby (Čísla v elipsách predstavujú veľkosti skupín získané po delení. Percentá pod vrcholmi predstavujú relatívny počet ľudí s lepším verzus horším výsledkom. SCLGSI = Symptom Checklist-90-R Global Severity Index). Obrázok modelu je prebraný a prepísaný z \cite[s. 454]{online-psychoterapy}}\label{fig:ClinicalTree}
\end{figure}

Vytvorený strom metódou CART obsahoval desať listových vrcholov a jeho chybovosť bola 18\%. Zo 448 pacientov, ktorým bol predpovedaný zlý výsledok malo 86\% skutočne zlý výsledok. Zo 182 pacientov s predikovaným dobrým výsledkom malo 74\% dobrý výsledok. Vytvorený strom je na obrázku \ref{fig:ClinicalTree}. Obrázok môže slúžiť ako učebnicová ukážka stromu vytvorená algoritmom CART.
\chapter{Genetické algoritmy}
Hlavným cieľom práce je optimalizovanie tvorby rozhodovacích stromov pomocou genetických algoritmov (GA). Preto 
v tejto kapitole predstavíme genetické algoritmy a podrobnejšie popíšeme jeho jednotlivé časti, ktoré budeme ďalej potrebovať v ďalšej kapitole.

Táto kapitola je delená do viacerých oddielov. V oddieli \ref{kap2:2.1:Info} popíšeme základné informácie o genetických algoritmoch. Popis rôznych kódovaní jedincov je v oddieli \ref{kap2:2.2:Coding}. Rôzne druhy funkcií, používaných na vyhodnotenie jedincov, zadefinujeme v oddieli \ref{kap2:2.3:Fitnesses}. Ďalej v časti \ref{kap2:2.4:Selection} predvedieme najzákladnejšie metódy na výber jedincov. Genetické operátory používané v genetických algoritmoch zavedieme v oddieli \ref{kap2:2.5:Operators}. Nakoniec v časti \ref{kap2:2.6:Parameters} popíšeme problém výberu správnych parametrov v genetických algoritmoch.

\section{Základný popis genetických algoritmov}\label{kap2:2.1:Info}
Genetické algoritmy sú optimalizačné techniky a prehľadávacie heuristiky, ktoré napodobujú princíp prirodzeného výberu vychádzajúceho z darwinovskej teórie. V tomto oddieli ponúkneme úvod ku genetickým algoritmom s popisom, ako by mal taký genetický algoritmus vyzerať. Informácie o genetických algoritmoch sme čerpali z knihy \cite{kap2-evolution} a úvod do oblasti genetiky zase z voľne dostupnej učebnice biológie \cite{online-biology}. Ďalšie doplňujúce poznatky boli získané z online zdrojov \cite{wiki-evolution,wiki-genetics}.

V časti \ref{kap2:2.1:2.1.1:Genetics} najprv z pohľadu biológie vysvetlíme základné pojmy z genetiky. Ďalej v časti \ref{kap2:2.1:2.1.2:AboutGeneticAlgo} popíšeme ako pracujú genetické algoritmy. Nakoniec v časti \ref{kap2:2.1:2.1.3:SimpleGeneticAlgo} predvedieme zjednodušený GA, ktorého hlavná myšlienka je rovnaká aj pri komplikovaných genetických algoritmoch.

\subsection{Genetika z pohľadu Biológie}\label{kap2:2.1:2.1.1:Genetics}
Na popísanie a zadefinovanie ďalších súčastí genetických algoritmov je dobré pripomenúť niektoré základné pojmy z oblasti genetiky. 

Bunka je základnou stavebnou jednotkou každého organizmu. Každá bunka má uložený popis toho, ako sa má správať pri typických bunečných procesoch. Tento popis sa nazýva genetická informácia a v organizmoch je uskladnená v DNA, čo je obrovská molekula, polymér zložený z veľkého množstva pospájaných nukleotidov (kyselina fosforečná s deoxyribózou a naviazanou dusíkatou bázou). Tieto nukleotidy a konkrétne ich báze sú zodpovedné zato, ako vyzerá genetický kód tejto genetickej informácie. Malý kúsok vybraný zo sekvencie DNA, ktorý kóduje nejakú funkčnosť bunky alebo proces v bunke, sa nazýva gén. Presnejšie vyjadrenie by bolo, že gén je zodpovedný za tvorbu proteínov, ktoré sú ovplyvňujú funkčnosť bunky. Gén si je teda možné predstaviť ako významnú črtu organizmu (napríklad farba očí). Rôzne nastavenia génu nazývame alely (modré, hnedé, zelené oči).

DNA je kvôli jej veľkosti ďalej komprimovaná do štruktúry nazývaná chromozóm. Chromozóm sa nachádza v jadre bunky. Organizmy môžu mať v jadre viacero rôznych chromozómov. Podľa typu organizmu môžu byť tieto chromozómy spárované (diploidné organizmy) alebo nespárované (haploidné organizmy). Genóm organizmu potom nazývame genetický materiál získaný zo všetkých chromozómov dokopy. Určitá množina génov z genómu organizmu sa nazýva genotyp. Vo vývojovom štádiu jedinca ovplyvňuje genotyp výzor a ďalšie charakteristiky organizmu (výška, farba očí, IQ, ...) nazývané fenotyp.

\begin{figure}[h]
\centering
\centerline{\mbox{\includegraphics[width=250pt]{../img/kap2/DNA.pdf}}}
\caption{Na obrázku je príklad bunky s jadrom, v ktorom sú viditeľné chromozómy. Jeden z chromozómov je v spodnej časti roztiahnutý do DNA, ktorej jednotlivé časti sú popísané. Značky A,G,T predstavujú nukleotidové báze adenín, guanín a cytozín.}\label{fig:DNA}
\end{figure}

Úlohou každého živého organizmu je rozšíriť svoju genetickú informáciu a tým pádom aj jeho gény. Sexuálna reprodukcia dvoch rodičovských organizmov je komplikovaný proces (zdroj \cite{online-shuffling} podáva rozumné vysvetlenie celého procesu). Konkrétne nás budú zaujímať dve štádia reprodukcie, ktoré funguje na úrovni chromozómov.
\begin{itemize}
\item Kríženie jednoduchých buniek alebo komplexných živočíchov, kde každá dvojica chromozómov si medzi sebou vymieňa gény a vytvára nové chromozómy.
\item Mutácia vytvorených chromozómov, pri ktorej sú niektoré nukleotidy pozmenené kvôli chýbam pri krížení.
\end{itemize}
Po reprodukcii dostávame nové chromozómy, ktoré tvoria genóm novo vytvoreného organizmu. Na obrázku \ref{fig:Chromosomes} je príklad reprodukcie nespárovaných chromozómov pre haploidné organizmy.

\begin{figure}[h]
\centering
\centerline{\mbox{\includegraphics[width=350pt]{../img/kap2/chromosomes.pdf}}}
\caption{Príklad zjednodušeného procesu reprodukcie chromozómov pre nového jedinca. Na obrázku a) sú chromozómy z rodičov. Gény v chromozómoch na obrázku b) sú vyznačené čiernymi plnými krúžkami. Kríženie génov v chromozómoch (vyznačené šípkami) a mutácia jedného z génov (šedý krúžok) je na obrázku c). Na obrázku d) sú vzniknuté chromozómy nového jedinca.}\label{fig:Chromosomes}
\end{figure}

Z biologického hľadiska je nutné popísať kvalitu nových organizmov. Tú je možné vyhodnotiť pomocou pravdepodobnosti, že sa jedinec dožije ďalšieho párenia. Ďalším kritériom môže byť plodnosť potomka.

\subsection{Popis genetických algoritmov}\label{kap2:2.1:2.1.2:AboutGeneticAlgo}
V tejto časti zadefinujeme, čo sú to genetické algoritmy, pričom jednotlivé pojmy budú podobné tým, ktoré sme predstavili v predchádzajúcej časti. 

Genetické algoritmy patria medzi metaheuristiky, ktorých úlohou je prechádzať priestor riešení a z nich nájsť to najlepšie. Algoritmus pracuje s jedincami, na ktorých aplikujeme operátory kríženia a mutácie. Podľa kvality jedinca rozhodujeme, či sa bude s jedincom ďalej pracovať. V genetických algoritmoch získavame riešenia z jedincov po príslušnom dekódovaní. Jedinec a chromozóm sú v rámci genetických algoritmov synonymami. Kvalita vytvorených jedincov záleží na tom, ako dobre tento jedinec rieši po dekódovaní zadanú úlohu. Väčšina implementácií pracuje s jedincami, ktoré by sa dali popísať ako haploidné (nespárované chromozómy).

Tieto algoritmy fungujú v cykloch, ktoré predstavujú generácie. V každom cykle pracujeme s množinou jedincov, populáciou pre tú danú generáciu. V každej generácii sa snažíme zlepšiť aktuálne riešenie vytváraním nových jedincov. V konečnej generácii je vytvorená taká populácia, ktorá obsahuje tých najvhodnejších jedincov, medzi ktorými sa nachádza ten najlepší.

Každý problém môže mať odlišnú reprezentáciu jedincov. S týmto súvisí úloha, ako rozumne a správne zakódovať jedincov. Takýchto kódovaní je mnoho, no medzi najznámejšie patrí binárne kódovanie, znakové kódovanie, kódovanie reálnymi číslami, ale aj kódovanie pomocou stromov, gramatík a mnoho ďalších. Napríklad pri použití binárneho kódovania sú jedinci reprezentovaní ako bitové reťazce. V tomto prípade sú jednotlivé prvky alebo skupiny prvkov v tomto poli génmi jedinca. Alela je v prípade bitovej hodnoty 0 alebo 1. Pri iných reprezentáciách sú alely komplikovanejšie.

Kríženie v genetických algoritmoch je realizované výmenou hodnôt medzi rodičmi. To o aké hodnoty sa jedná záleží na zvolenej reprezentácií (hodnoty v reťazci, stromy, ...). Typ kríženia, ktorý môžeme použiť, závisí od zvoleného kódovania a od jeho reprezentácie jedincov. 
Mutácia je realizovaná zmenou hodnoty už v kríženom jedincovi, pričom podobne ako pri krížení je závislá na kódovaní jedincov a od použitej reprezentácie. 

Napríklad v prípade kódovania, pri ktorom sú jedinci reprezentovaní reťazcom je typ kríženia obmedzený na bodové a uniformné kríženie. Rozumnou mutáciou je v tomto prípade invertovanie bitu.

\subsection{Jednoduchý genetický algoritmus}\label{kap2:2.1:2.1.3:SimpleGeneticAlgo}
Podľa doposiaľ zadefinovaných pojmov môžeme predviesť, ako taký genetický algoritmus pracuje. Majme daný problém P, ktorý chceme vyriešiť. Pre zjednodušenie uvažujme jedincov s binárnym kódovaním. Zjednodušený Algoritmus \ref{fig:geneticAlgoritm} potom popisuje jednotlivé kroky genetického algoritmu, ktorý rieši problém P.
\begin{algorithm}
\floatname{algorithm}{Algoritmus}
\caption{Kroky zjednodušeného genetického algoritmu, ktorý pracuje s binárne kódovanými jedincami.}\label{fig:geneticAlgoritm}
$A$ - Aktuálna populácia \\
$O$ - Potomkovia \\
$k$ - Dĺžka reťazca/chromozómu/riešenia \\
$fP$ - Kritérium na vyhodnotenie jedincov pre problém $P$\\
$p_{k}$ - Pravdepodobnosť kríženia \\
$p_{m}$ - Pravdepodobnosť mutácie \\
$max\_gen$ - Maximálny počet generácií \\
$max\_pop$ - Maximálny počet jedincov v populácií \\
\begin{algorithmic}

\State \parbox[t]{400pt}{$A \gets $ náhodne vygenerovaná populácia s $n$ jedincami (každé z nich dĺžky $k$).}

\For{$i = 1$ kým $i \leq max\_gen$}
	\State Krok()
\EndFor
\\
\Procedure{Krok}{\mbox{}}
	\State $O \gets \{\}$
	\State Spočítaj kvalitu každého jedinca pomocou kritéria $fP$.	
	\While{$\lvert O \rvert < max\_pop$}
		\State \parbox[t]{350pt}{$J1,J2 \gets $ dva jedinci z populácie $A$ (výber s opakovaním). Pravdepodobnosť že vyberieme jedinca je úmerná hodnote kritéria $fP$.}
		\State \parbox[t]{350pt}{$O1,O2 \gets $ \textsc{Kríženie(J1,J2)}}
		\State \parbox[t]{350pt}{$O1,O2 \gets $ \textsc{Mutácia(O1,O2)}}
		\State $O \gets \{O1,O2\} \cup O$.
	\EndWhile
	\State $A \gets O$.
\EndProcedure
\\
\Function{Kríženie}{$J1$,$J2$} \\
\Return{\parbox[t]{350pt}{$O1, O2$, ktorí vznikli prekrížením rodičov $J1, J2$ s pravdepodobnosťou $p_{k}$ v ľubovoľnom mieste reťazca. Bez použitého kríženia bude $O1 \gets J1$ a $O2 \gets J2$.}}
\EndFunction
\\
\Function{Mutácia}{$J1$,$J2$} \\
\Return{\parbox[t]{350pt}{$O1, O2$, ktorí vznikli mutáciou jedincov $J1, J2$ v jednotlivých hodnotách reťazca. Mutácia je na hodnote v reťazci použitá s pravdepodobnosťou $p_{m}$.}}
\EndFunction
\end{algorithmic}
\end{algorithm}
Každá iterácia algoritmu predstavuje novú generáciu jedincov. Počet generácii závisí na typu úlohu, no zvyčajne ich býva 50 až 500. V poslednom behu algoritmu dostávame konečnú populáciu, v ktorej je viacero kvalitných jedincov. Na zreprodukovanie výsledkov je kvôli stochastickému správaniu algoritmu nutné zvoliť rovnaký seed (pri používaní pseudo náhodných generátorov). 

Väčšina genetických algoritmov pracuje podobne ako zjednodušený algoritmus. Najväčšia zmena býva vtom akú taktiku na výber jedincov a aké kríženie a mutáciu použijeme. Komplikovanejšie genetické algoritmy taktiež používajú princíp elitizmu, ktorý ovplyvňuje aké percento rodičovských jedincov sa dostane do ďalšej generácie.

\section{Kódovanie jedincov}\label{kap2:2.2:Coding}
V predchádzajúcom oddieli sme uviedli, že jedince môžu byť rôzne kódované. Výber správneho kódovania záleží hlavne na typu úlohy, ktorú riešime. Každé jedno z kódovaní pracuje s určitou reprezentáciou jedinca. Aplikácie genetických algoritmov používajú buď reprezentácie s pevne zadanou dĺžkou, alebo s variabilnou dĺžkou. Existuje mnoho kódovaní jedincov, pričom tie najznámejšie sú
\begin{itemize}
\item binárne, pri ktorom reprezentujeme jedinca bitovým reťazcom/poľom určitej dĺžky,
\item znakové, kde je jedinec reprezentovaný reťazcom, ktorého hodnoty sú z konečnej množiny symbolov,
\item celočíselné, kde je jedinec reprezentovaný reťazcom celočíselných hodnôt,
\item kódovanie s reálnymi číslami, pri ktorom reprezentujeme jedinca reťazcom, ktorého hodnoty sú reálne čísla,
\item stromové, kde sú jedince reprezentované ako stromy z teórie grafov,
\item a mnoho ďalších ako je permutačné kódovanie, maticové kódovanie, kódovanie založené na gramatikách, ...
\end{itemize} 

V tomto oddieli budú jednotlivé časti popisovať rôzne kódovania jedincov. V časti \ref{kap2:2.2:2.2.1:Binary} popíšeme binárne kódovanie, ktoré bolo historicky prvým kódovaním jedincov v genetických algoritmoch. Kódovanie s reálnymi číslami zadefinujeme v časti \ref{kap2:2.2:2.2.2:RealValued}. Nakoniec v časti \ref{kap2:2.2:2.2.3:Tree} uvedieme kódovanie pomocou stromov z teórie grafov.
\subsection{Binárne kódovanie}\label{kap2:2.2:2.2.1:Binary}
Binárne kódovanie je prvý typ kódovania, s ktorým pracovali prvé, originálne genetické algoritmy.
Tento typ kódovania je najčastejšie používaným kódovaním jedincov už od dôb Hollanda. Kódovanie uvažuje jedincov, ktorí sú reprezentovaní ako reťazec, ktorého hodnoty sú binárne (obrázok \ref{fig:Encoding}a). Obľúbenosť si získalo hlavne kvôli tomu, že 
\begin{enumerate}
\item Holland a jeho študenti sa venovali tomuto kódovaniu najdlhšie, čo malo za následok, že väčšinu teoretických poznatkov máme dokázaných pre binárne kódovanie (príkladom môže byť veta o schémach).
\item Pre binárne kódovanie existujú určité heuristiky, ako správne zvoliť parametre genetického algoritmu.
\end{enumerate}
Veľa teoretických poznatkov sa dá rozšíriť aj na nebinárne kódovanie, no tieto rozšírenia nie sú tak dobre zdokumentované, ako to je v prípade originálneho GA.

Praktické použitie genetického algoritmu, ktorý pracuje s binárne kódovanými jedincami je napríklad problém batohu. Batoh si je možné predstaviť ako binárny reťazec. Počet predmetov, ktoré je možné vložiť do batohu je $N$. Jednička v reťazci na mieste $i$ predstavuje vloženie predmetu $i$ do batohu, pričom batoh nezostane zároveň preplnený.
\subsection{Kódovanie reálnymi číslami}\label{kap2:2.2:2.2.2:RealValued}
Historicky prvé pokusy boli postavené na zakódovaní reálnych čísiel do bitovej reprezentácie. Dnes sa tento typ zakódovania nepoužíva dokým nie nato je naozaj dobrý dôvod. Neskoršie bolo kódovanie reálnych čísiel postavené na prirodzenejšej reprezentácii, ktorá zostala až dodnes. Takáto reprezentácia využíva reťazce, ktorého hodnoty sú čísla s plávajúcou čiarkou (floatové čísla).
Príklad jedinca kódovaného reálnymi číslami je na obrázku \ref{fig:Encoding}b. Genetické operátory musia byť zvolené tak, aby rešpektovali takéto reťazce reálnych čísiel.

Príkladom použitia genetického algoritmu, ktorý pracuje s jedincami kódovanými reálnymi číslami je mnoho kriteriálna optimalizácia funkcií. Pri tejto úlohe sa snažíme minimalizovať (alebo maximalizovať) hodnotu viacero zadaných funkcií, čo je ale NP ťažké. Reálne kódovaný jedinec predstavuje bod v $n$-dimenzionálnom priestore. Genetický algoritmus prehľadáva priestor riešení a na konci vráti takého jedinca, ktorý čo najlepšie minimalizuje všetky funkcie.
\begin{figure}[h]
\centering
\centerline{\mbox{\includegraphics{../img/kap2/encoding.pdf}}}
\caption{Príklady kódovania jedincov. Na obrázku a) je binárne kódovaný jedinec/chromozóm. Na obrázku b) je jedinec kódovaný reálnymi číslami. Obidva chromozómy sú reprezentovaní reťazcami.}\label{fig:Encoding}
\end{figure}

\subsection{Kódovanie stromami}\label{kap2:2.2:2.2.3:Tree}
Pri netypických úlohách, ktorých jedincov by bolo ťažké reprezentovať pomocou reťazcov, bolo nutné zaviesť také kódovania, ktoré by zrozumiteľnejšie popisovalo jedincov. Pri úlohách, ktorých riešenia sa dajú jednoducho previesť na strom bolo rozumné použiť na kódovanie jedincov túto stromovú štruktúru. Kvôli tejto štruktúre bolo nutné zaviesť rôzne nové mutácie a kríženia, ktoré by dávali zmysel pre toto kódovanie. 

Existuje veľa problémov v genetických algoritmoch, ktoré pracujú s jedincami kódovanými pomocou stromov. Stromy v GA majú tendenciu narastať. V biologickej evolúcii existuje určitý fyzikálny limit, ktorý bráni nadmernému rastu. Ten je v GA možné dosiahnuť reštrikciou veľkosti stromu alebo obmedzením maximálnej hĺbky. To je možné dosiahnuť buď pomocou kritérií kvality stromov alebo pomocou operátorov, ktoré nezväčšujú jedincov (prípadne ich zmenšujú).

Príklad GA pracujúceho s jedincami kódovanými stromami je metóda genetického programovania (GP). Hlavnou úlohou genetického programovania je hľadať programy (naši jedinci), ktoré riešia zadané problémy. Táto metóda sa v podstate skladá z množiny inštrukcií a kritéria na vyhodnotenie vytvoreného programu. Stromy reprezentujúce programy majú vo vnútorných vrcholoch operátory (obyčajne aj zložité) a v listoch zase konštanty. Príklad takéhoto stromového jedinca je na obrázku \ref{fig:TreEncoding}.
\begin{figure}[h]
\centering
\centerline{\mbox{\includegraphics{../img/kap2/treencoding.pdf}}}
\caption{Príklad kódovania jedincov pomocou stromov. Metóda genetického programovania pracuje s rovnako vyzerajúcimi jedincami. Strom na tomto obrázku kóduje funkciu $\left(X - \dfrac{Y}{8}\right) + \left(7 * \cos(Z)\right)$. Vnútorné vrcholy predstavujú operácie a listy predstavujú konštanty. }\label{fig:TreEncoding}
\end{figure}

\section{Vyhodnotenie jedincov}\label{kap2:2.3:Fitnesses}
Dôležitou častou genetických algoritmov je nejakým spôsobom ohodnotiť kvalitu jedincov. V GA určuje kvalita jedinca väčšinou kvalitu riešenia zadanej úlohy, no zároveň popisuje aj to, ako často sa jedinec zúčastní procesu reprodukcie. Kritérium na ohodnotenie kvality jedinca je v GA nazývaná fitness. Úlohy, ktoré hodnotia jedinca iba podľa jedného kritéria nazývame jedno kriteriálne (angl. SOEA). 

V prípade úloh, pri ktorých určujeme kvalitu jedinca pomocou dvoch a viacerých kritérií (mnoho kriteriálne problémy angl. MOEA) je potrebné pozmeniť genetický algoritmus tak, aby s nimi vedel nejako rozumne pracovať. To je možné dosiahnuť rôznymi metódami.
V časti \ref{kap2:2.3:2.3.1:Weighted} predvedieme metódu, ktorá rieši problém viacerých kritérií agregáciou fitness. Ďalej v časti \ref{kap2:2.3:2.3.2:Priority} popíšeme metódu ktorá určuje kvalitu jedincov podľa kritérií s danou prioritou. Nakoniec v časti \ref{kap2:2.3:2.3.3:Pareto} zadefinujeme, ako sa dá určiť kvalita jedincov (pomocou viacerých kritérií) podľa paretovej fronty.

\subsection{Agregácia fitness}\label{kap2:2.3:2.3.1:Weighted}
Podstatou tejto techniky je vytvoriť zo zadaných kritérií $f_{i}$, $i=1,\ldots,k$ jednu hodnotu $f$. V podstate sa snažíme previesť mnoho kriteriálny problém na jedno kriteriálny problém. To sa dá docieliť spočítaním váženej sumy jednotlivých kritérií, kde váhy kritérií sú predom definované vo váhovom vektore $\mathbf{w}$ veľkosti $k$.
\begin{equation}
f = \sum_{i=1}^{k} w_{i} * f_{i} \nonumber
\end{equation}
Problém tejto metódy spočíva v správnom nastavení váhového vektora $\mathbf{w}$.
\subsection{Prioritná fitness}\label{kap2:2.3:2.3.2:Priority}
Prioritná fitness rieši mnoho kriteriálny problém tak, že priraďuje priority ku zadaním kritériám $f_{i}$, $i=1,\ldots,k$. Priorita môže byť predom určená podľa poradia v akom sú kritéria zadané (prvé kritérium má najväčšiu prioritu). Postup pri porovnaní dvoch jedincov potom funguje následovne:
\begin{enumerate}
\item Pre oboch jedincov spočítame hodnoty kritéria s najväčšou prioritou.
\item Keď sú spočítané hodnoty rôzne, tak budú jedinci usporiadaní podľa nich.
\item Keď sú spočítané hodnoty rovnaké, tak ďalej porovnávame jedincov podľa ďalšieho kritéria v poradí.
\end{enumerate}

\subsection{Pareto fitness}\label{kap2:2.3:2.3.3:Pareto}
Táto metóda rieši mnoho kriteriálne problémy pomocou paretovej fronty. Na pochopenie pareto fronty je najprv nutné zadefinovať niektoré pojmy.
\begin{def-sk}
Nech $F = \{f_{1}, \ldots, f_{k}\}$ je množina kritérií (fitness funkcií) podľa ktorých chceme ohodnotiť jedincov. Ďalej nech $x$ a $y$ sú dva jedince. Jedinec $x$ slabo dominuje jedinca $y$, práve vtedy keď $f_{i}(x) \leq f_{i}(y)$, pre $i = 1,\ldots,k$.
\end{def-sk}
\begin{def-sk}
Nech $F = \{f_{1}, \ldots, f_{k}\}$ je množina kritérií (fitness funkcií) podľa ktorých chceme ohodnotiť jedincov. Ďalej nech $x$ a $y$ sú dva jedince. Jedinec $x$ dominuje jedinca $y$, práve vtedy keď ho slabo dominuje a existuje $j$: $f_{j}(x) < f_{j}(y)$.
\end{def-sk}
\begin{def-sk}
Nech $x$ a $y$ sú dva jedinci. Potom hovoríme, že $x$ a $y$ sú neporovnateľní, keď $x$ nedominuje $y$ a $y$ nedominuje $x$.
\end{def-sk}
\begin{def-sk}
Nech $x$ a $y$ sú dva jedinci. Jedinec $x$ nedominuje $y$ práve vtedy, keď $y$ slabo dominuje $x$ alebo sú $x$ a $y$ neporovnateľní.
\end{def-sk}
\begin{def-sk}
Paretova fronta je definovaná ako množina jedincov, ktorí nie sú dominovaní žiadnym iným jedincom (obrázok \ref{fig:Pareto}).
\end{def-sk}
Potom jedným z postupov ako vybrať správneho jedinca z paretovej fronty je postavený na tom, aké kritérium chceme najviac optimalizovať. Existujú aj iné techniky, ktoré sú založené na dominancii jedincov. Najznámejšie z nich sú NSGA a NSGAII \cite{online-NSGA}.

\begin{figure}[h]
\centering
\centerline{\mbox{\includegraphics{../img/kap2/pareto.pdf}}}
\caption{Príklad ako vyzerá paretova fronta pre dve kritéria $f_{1}$ a $f_{2}$. Jedinci sú označení čiernymi štvorcami. Na paretovej fronte (čiara spájajúca štvorce) sú jedinci, ktorí nie sú dominovaní žiadnym iným jedincom.}\label{fig:Pareto}
\end{figure}

\section{Výber jedincov}\label{kap2:2.4:Selection}
Po zadefinovaní kódovania a kritérií na vyhodnotenie jedincov môžeme definovať metódy na výber jedincov k reprodukcii, ďalej ako selekcia (angl. selection methods). Hlavnou úlohou týchto metód je vyberať tých najkvalitnejších jedincov z aktuálnej generácie (podobnosť s prirodzeným výberom z prírody) na vytvorenie novej, kvalitnejšej generácie potomkov.

Pri riešení úloh je nutné rozumne zvoliť metódy výberu jedincov z populácie. Pri zvolení selekcie, ktorá by vyberala iba tých najlepších jedincov z populácie vedie ku novej generácii neoptimálnych jedincov. Toto by malo za následok zníženie diverzity, čo by viedlo k rýchlej konvergencii genetického algoritmu do nejakého lokálneho optima. Naopak výber selekcie, ktorá by vyberala jedincov náhodne vedie k príliš dlhej konvergencii genetického algoritmu. Univerzálny návod nato, ako zvoliť vhodnú selekciu v GA nie je doteraz známy.

Príklady metód na výber jedincov sú: ruletová selekcia, turnajová selekcia, SUS (angl. stochastic universal sampling), sigmoidálna selekcia, \ldots. V tomto oddieli popíšeme tie najznámejšie z nich. V prvej časti \ref{kap2:2.4:2.4.1:RoulletteWheel} predvedieme ako funguje historicky prvá, ruletová selekcia, ktorá bola použitá v originálnom GA. V časti \ref{kap2:2.4:2.4.2:Tournament} predstavíme ako funguje výkonnejšia turnajová selekcia. V poslednej časti \ref{kap2:2.4:2.4.3:Elitism} zadefinujeme pojem elitizmus a akú úlohu zohráva pri výbere jedincov.

\subsection{Ruletová selekcia}\label{kap2:2.4:2.4.1:RoulletteWheel}
Originálny GA od Hollanda používal selekciu, ktorá vyberala jedincov podľa ich kvality. Takýto typ selekcie sa dá implementovať pomocou "rulety", ktorá funguje tak, že jedincom z populácie priraďujeme jednotlivé kúsky tejto pomyselnej "rulety". Veľkosť priradeného kúsku závisí na fitness hodnote jedinca. Pravdepodobnosť výberu je potom úmerná tejto fitness.

Vytvorenie rulety prebieha v dvoch krokoch:
\begin{enumerate}
\item Sčítaj fitness každého jedinca.
\item Prepočítaj pravdepodobnosti výberu jedincov (fitness jedinca vydelená sumou).
\end{enumerate}
Na vytvorenie $n$ jedincov otáčame touto "ruletou" $n$-krát. Vybraný jedinec je zaradený do novej populácie potomkov, na ktorú sú ďalej použité rôzne genetické operátory.
\subsection{Turnajová selekcia}\label{kap2:2.4:2.4.2:Tournament}
Turnajová selekcia je výpočtovo efektívnejšou a paralelne prístupnejšia metóda na výber jedincov. Oproti ruletovej selekcii je efektívnejšia v tom, že nemusí dva razy prechádzať populáciou a výber jednotlivcov je celkovo zjednodušený. Turnajová selekcia s parametrom $k$ funguje v dvoch krokoch
\begin{enumerate}
\item Vyber dvoch náhodných jedincov z populácie.
\item Zvoľ náhodné číslo $r$, $0 < r < 1$. Keď je $r < k$, tak zvolíme jedinca s lepšou fitness. Inak zvolíme jedinca s menšou fitness. 
\end{enumerate} 

\subsection{Elitizmus}\label{kap2:2.4:2.4.3:Elitism}
Elitizmus o hodnote $k$, ktorý bol po prvý krát predstavený pánom De Jong, je špeciálnym typom výberu jedincov. Na vybraných jedincoch nie sú aplikované žiadne operátory. Namiesto toho ich priamo zaraďujeme do ďalšej generácie. Týmto donútime genetický algoritmus, aby si ponechal určité percento $k$ najlepších jedincov a tým neriskujeme ich stratu (krížením alebo mutáciou). Bolo ukázané, že použitie elitizmu o hodnote $k$ signifikantne zlepšuje efektivitu genetických algoritmov.

\section{Genetické operátory}\label{kap2:2.5:Operators}
V prípade kódovania, ktoré využíva reprezentáciu reťazcom môžeme použiť
\begin{itemize}
\item bodové kríženie, ktoré je založené na zvolení náhodných $n$ bodov. Takto zvolené body rozdelia reťazec na $n+1$ častí, pričom sa odpovedajúce časti vymenia s určitou pravdepodobnosťou. Tento typ kríženia sa ďalej pomenováva podľa počtu vybraných bodov (jednobodové, dvojbodové, ...).
\item uniformné kríženie, ktoré vypĺňa jednotlivé hodnoty z určitého rodiča. U každej položky hádžeme mincou, ktorá rozhoduje z ktorého rodiča hodnotu vyberieme.
\end{itemize}

\section{Nastavenie správnych parametrov}\label{kap2:2.6:Parameters}
Poslednou úlohou pri práci s genetickými algoritmami je problém nastavenia vhodných parametrov ako je veľkosť populácie, pravdepodobnosť kríženia, pravdepodobnosť mutácie prípadne hodnota elitizmu. Hľadanie riešenia pomocou GA závisí (pomimo vhodných operátorov) hlavne na týchto parametroch.

Naivným spôsobom by bolo optimalizovať parametre postupne. Toto v prípade GA nie je možné, keďže sa jednotlivé parametre navzájom ovplyvňujú. Problém správneho nastavenia parametrov či adaptácie parametrov je obľúbenou témou genetických algoritmov a existuje mnoho článkov a výskumov, ktoré sa tejto téme venujú. Napriek tomu neexistuje žiadny univerzálny návod a väčšinou sa používajú také hodnoty parametrov, ktoré už v minulosti dávali dobré výsledky. Pravdaže nastavenie parametrov, vo väčšine prípadov, závisí na typu úlohy. To znamená, že jedno špeciálne nastavenie, ktoré dobre fungovalo na jednom type úlohy môže dávať katastrofické výsledky na inom type úlohy.

Jeden z pánov, ktorý experimentoval s nastavením parametrov pričom dávali dobré výsledky, bol De Jong. Ten skúmal úlohy s binárnym kódovaním. Jeho experimenty uvádzali veľkosť populácie približne 50 až 100 jedincov. Pravdepodobnosť jednobodového kríženia bola približne 0.6 a pravdepodobnosť mutácie bola 0.001 na jeden bit. Iná skupina ľudí uviedla, že najvhodnejšími parametrami bola veľkosť populácie 20-30, pravdepodobnosť kríženia 0.75 až 0.95 a pravdepodobnosť mutácie na jeden bit 0.005 až 0.001. Jednou z možnosti ako nastaviť vhodné parametre je využiť ďalší GA na ich optimalizáciu.
\chapter{Genetická tvorba rozhodovacích stromov}\label{kap3:DTGA}
V časti \ref{kap1:2.7:DTTechniques} sme popísali vytváranie rozhodovacích stromov pomocou indukčných algoritmov. Vieme, že stromy vytvorený týmito algoritmami bývajú kvalitné, ale nemusia byť optimálne. To je kvôli tomu, že tieto algoritmy patria medzi tzv. hladné techniky. Ďalším spôsobom, ako vytvárať stromy je použiť techniku, ktorá prehľadáva priestor všetkých riešení (napríklad prehľadávaciu heuristiku). Pri použití niektorej z týchto techník môžeme dostať neoptimálne riešenie, ktoré môže byť dokonca horšie, než to vytvorené indukčným algoritmom. Vhodne zvolená a nastavená prehľadávacia heuristika ale môže nájsť riešenie, ktoré je významne lepšie než to vytvorené indukčným algoritmom. V kapitole \ref{kap2:GA} sme predviedli jednu z prehľadávacích heuristík nazývanú genetické algoritmy.

V tejto kapitole teda popíšeme genetický algoritmus, ktorý sa dá použiť na tvorbu rozhodovacích stromov a ktorý sme implementovali v nástroji Weka. Popíšeme, ako kódovať jedincov v populácií genetického algoritmu tak, aby títo jedinci čo najlepšie popisovali rozhodovacie stromy. Operátory zavedené v predchádzajúcej kapitole rozšírime o nové genetické operátory špeciálne navrhnuté na tvorbu rozhodovacích stromov. Ostatné súčasti genetického algoritmu iba pripomenieme a spojíme ich s pojmami z rozhodovacích stromov.

Vytváranie stromov genetickými algoritmami nie je novou technikou. Existujú práce, ktoré sa venujú tejto problematike, no tie väčšinou vytvárajú stromy od úplného začiatku. Niektoré z týchto prác spomenieme v kapitole \ref{kap6:SimilarWorks}. Taktiež treba uviesť, že voľne dostupných implementácií je málo a zatiaľ sa zdá, že neexistuje žiadna do nástroja Weka.

Najskôr v oddieli \ref{kap3:3.1:Intro} popíšeme úlohu tvorby stromov genetickými algoritmami. Ďalej v oddieli \ref{kap3:3.2:Encoding} pripomenieme kódovanie jedincov pomocou stromov a ako sa dá toto kódovanie pozmeniť na reprezentovanie rozhodovacích stromov. V oddieli \ref{kap3:3.3:Fitness} uvedieme fitness funkcie používané v takomto genetickom algoritme. Na konci, v oddieli \ref{kap3:3.4:Operators}, pripomenieme operátor kríženia stromov a zadefinujeme dve nové mutácie stromov.
\section{Základný popis}\label{kap3:3.1:Intro}
V tomto oddieli popíšeme genetický algoritmus, ktorý sa dá použiť na vytvorenie rozhodovacích stromov. Techniky tvorby rozhodovacích stromov pomocou genetických algoritmov môžeme rozdeliť na dve skupiny:
\begin{enumerate}
\item vytváranie stromov úplne od začiatku,
\item optimalizovanie už vytvorených stromov a zlepšenie ich prediktívnych schopností,
\end{enumerate}
pričom obidve skupiny majú svoje výhody aj nevýhody.

Vytváranie stromov pomocou genetických algoritmov od začiatku je pomalé a nájdenie vhodného riešenia môže trvať príliš dlho. Výhodou je, že stromy vytvárame od začiatku a teda nie sú ovplyvňované takými stromami, ktorými by sme rýchlo skonvergovali k lokálnemu optimu. 

Pri optimalizovaní vytvoreného stromu môžeme využiť to, že počiatočná generácia je už na počiatku kvalitná. V niektorých prípadoch by stačila iba malá zmena na týchto kvalitných stromoch a generalizačné schopnosti pozmeneného stromu by sa významne zlepšili. V iných prípadoch, kedy genetický algoritmus nedokáže vylepšiť generalizačné schopnosti, môžeme vylepšovať iné vlastnosti stromu (napr. veľkosť stromu, výšku stromu) pričom nezhoršíme generalizačné schopnosti stromu. Na úkor dodatočného času teda môžeme dostať rozhodovacie stromy, ktoré svojou kvalitou prevyšujú alebo sa aspoň vyrovnávajú tým vytvoreným indukčnými algoritmami, pričom zlepšujeme ostatné kritériá. Pre predstavu uvádzame dva obrázky rozhodovacích stromov vytvorených pre tú istú úlohu. Úloha pozostávala z klasifikácie náhodných dát, vygenerovaných generátorom \verb|RDG1|, ktorý ponúka nástroj Weka. Prvý zo stromov bol vytvorený indukčným algoritmom (Obrázok \ref{fig:treeIA}) a druhý bol optimalizovaný genetickým algoritmom (Obrázok \ref{fig:treeGA}). Ani jeden z týchto modelov nemusí byť významne lepší v predikovaní, no veľkosť stromu optimalizovaného genetickým algoritmom je výrazne menšia než veľkosť stromu vytvoreného indukčným algoritmom.

\begin{figure}[h]
\centering
\centerline{\mbox{\includegraphics[width=400pt]{../img/kap3/tree1.pdf}}}
\caption{Príklad stromu vytvoreného nejakým indukčným algoritmom.}\label{fig:treeIA}
\end{figure}

V tejto práci sa kvôli jeho výhodám zameriame na optimalizovanie vytvorených stromov genetickými algoritmami. Pritom počiatočné stromy nebudú náhodné, ale do určitej miery kvalitné. Kvalitu stromu vynútime použitím niektorého z algoritmov spomínaných v časti \ref{kap1:2.6:DTEvaluation}. 

Celým procesom optimalizácie stromov sa snažíme čo najviac zlepšiť (tj. zmenšiť) generalizačnú chybu, pričom môžeme vylepšovať aj ďalšie kritéria stromu, ako napríklad veľkosť stromu. Z časti \ref{kap1:2.6:2.6.1:Generalize} vieme, že generalizačnú chybu môžeme aproximovať empirickou chybou na nezávislej testovacej množine. Dátovú množinu $D$ teda rozdelíme na trénovaciu podmnožinu $T$ a validačnú podmnožinu $V$. Na ešte lepšiu aproximáciu generalizačnej chyby môžeme použiť $n$-násobnú krížovú validáciu viz. časť \ref{kap1:2.6:2.6.1:Generalize}. Obidva prístupy sú založené na ohodnotení kvality vytvorených stromov. Algoritmus by mal kvôli tomu poskytovať nejakú formu aproximácie generalizačnej chyby.

\begin{figure}[h]
\centering
\centerline{\mbox{\includegraphics[width=300pt]{../img/kap3/tree2.pdf}}}
\caption{Príklad stromu vytvoreného indukčným algoritmom a optimalizovaný genetickým algoritmom.}\label{fig:treeGA}
\end{figure}

Počiatočné rozhodovacie stromy, ktoré budeme optimalizovať, teda budú vybudované z trénovacej množiny $T$. Operátory budú slúžiť na prehľadávanie priestoru riešení a vyhľadanie nových rozhodovacích stromov. Typy operátorov sú zvolené tak, aby čo najviac diverzifikovali populáciu. Na vyberanie jedincov je vhodné použiť elitizmus (kvôli udržaniu najlepšieho riešenia) a turnajovú selekciu hlavne kvôli jej výhodám spomínaným v časti \ref{kap2:2.4:2.4.2:Tournament}. Genetickým algoritmom vytvárame nové stromy, ktorých kvalitu určíme podľa zvolených kritérií. Použijeme výpočet fitness z jedného alebo viacerých kritérií viz. časť \ref{kap2:2.3:Fitnesses}. Validačná množina $V$ bude slúžiť na otestovanie modelu a na odhad generalizačnej chyby.

Pri tomto genetickom algoritme použijeme selekciu prostredia z časti \ref{kap2:2.4:2.4.4:EnvironmentalSelection}, čo spolu s elitizmom zabráni strate tých najlepších jedincov a zaručí, že populácia zostane dostatočne diverzifikovaná.

V časti \ref{kap2:2.1:2.1.2:AboutGeneticAlgo}, konkrétne na Obrázku \ref{fig:GAdiagram}, sme zaviedli diagram genetického algoritmu. Podľa tohoto diagramu sa riadi aj chovanie genetického algoritmu na optimalizovanie stromov.
V algoritme \ref{fig:DTGeneticAlgo} uvádzame pseudokód tohoto genetického optimalizovania stromov a popisujeme jeho jednotlivé súčasti.
Pseudokód je veľmi podobný tomu spomenutému v časti \ref{kap2:2.1:2.1.3:SimpleGeneticAlgo} v algoritme \ref{fig:geneticAlgoritm}.

\begin{algorithm}
\floatname{algorithm}{Algoritmus}
\caption{Kroky genetického algoritmu, ktorý optimalizuje rozhodovacie stromy vytvorené indukčným algoritmom.}\label{fig:DTGeneticAlgo}
$A$ - aktuálna populácia rozhodovacích stromov \\
$O$ - potomkovia \\
$e$ - hodnota elitizmu, $e \in [0,1]$ \\
$s$ - zvolená metóda selekcie (turnaj, ruleta) \\
$sp$ - zvolená metóda selekcie prostredia (turnaj, ruleta). \\
$f_1,\ldots,f_p$ - $p$ kritérií na vyhodnotenie rozhodovacích stromov pre problém $P$\\
$k_{1},\ldots,k_{q}$ - pravdepodobnosti $q$ operátorov kríženia\\
$m_{1},\ldots,m_{r}$ - pravdepodobnosti $r$ operátorov mutácie\\
$max\_gen$ - maximálny počet generácií \\
$max\_pop$ - maximálny počet jedincov v populácií \\
\bigskip
\begin{algorithmic}[1]
\State \parbox[t]{375pt}{$A \gets $ populácia $max\_pop$ stromov, vytvorených nejakým indukčným algoritmom.}
\For{$i = 1$ \textbf{to} $max\_gen$} 
	\State \textsc{Krok($A$)}
\EndFor
\\
\Procedure{Krok}{$A$}
	\State $O \gets \{\}$
	\State \textsc{VyhodnoťPopuláciu($A$)}
	\While{$\lvert O \rvert < max\_pop$}
		\State \parbox[t]{350pt}{$J_1,J_2 \gets $ \textsc{VyberJedincov($A$)}}
		\State \parbox[t]{350pt}{$K_1,K_2 \gets $ \textsc{Kríženie($J_1$,$J_2$)}}
		\State \parbox[t]{350pt}{$O_1 \gets $ \textsc{Mutácia($K_1$)}}
		\State \parbox[t]{350pt}{$O_2 \gets $ \textsc{Mutácia($K_2$)}}		
		\State $O \gets \{O_1,O_2\} \cup O$.
	\EndWhile
	\State $E \gets$ \textsc{Elitizmus($A$,$e$)}.
	\State $A \gets  E$ $\cup$ \textsc{VyberJedincovProstredia($O$)} 
\EndProcedure
\\
\Procedure{VyhodnoťPopuláciu}{$A$}
\State \parbox[t]{350pt}{každý rozhodovací strom $x$ z aktuálnej populácie $A$ ohodnoť kritériami $f_i, 1 \leq i \leq p$.}
\EndProcedure
\\
\Function{VyberJedincov}{$A$}
\State \Return \parbox[t]{300pt}{$O_1$,$O_2$ vybraných metódou selekcie $s$ z populácie $A$.}
\EndFunction
\\
\Function{Kríženie}{$J_1$,$J_2$}
\State $K_1, K_2 \gets J_1, J_2$
\For{$i = 1$ \textbf{to} $q$} 
\State $r \gets [0,1]$
\If{$r < \sum_{j=1}^{i}k_j$}
	\State \parbox[t]{300pt}{kríženie $i$ aplikuj na stromy $K_1, K_2$.}
\EndIf
\State \Return $K_1$,$K_2$
\EndFor 
\EndFunction
\algstore{myalg}
\end{algorithmic}
\end{algorithm}

\begin{algorithm}
\ContinuedFloat
\caption{pokračovanie...}
\begin{algorithmic}[1]
\algrestore{myalg}
\Function{Mutácia}{$K$}
\State $M \gets K$
\For{$i = 1$ \textbf{to} $r$} 
\State $r \gets [0,1]$
\If{$r < m_i$}
	\State \parbox[t]{300pt}{mutáciu $i$ aplikuj na strom $M$ (na jedinca môže byť použitých viacero mutácií).}
\EndIf
\EndFor 
\State \Return $M$
\EndFunction
\\
\Function{Elitizmus}{$A$,$e$}
\State \Return \parbox[t]{300pt}{$e |A|$ najlepších jedincov}
\EndFunction
\\
\Function{VyberJedincovProstredia}{$O$}
\State \Return \parbox[t]{300pt}{rozhodovacie stromy vybrané z $O$ selekciou prostredia $sp$. Počet vybraných stromov bude rovný $(1-e)max\_pop$.}
\EndFunction
\end{algorithmic}
\end{algorithm}

\section{Kódovanie jedincov}\label{kap3:3.2:Encoding}
Jedinci v tomto genetickom algoritme sú rozhodovacie stromy. Z časti \ref{kap1:2.3:2.3.2} vieme, že rozhodovací strom ma štruktúru stromu. Najrozumnejším kódovaním je preto kódovanie pomocou stromov, ktoré bolo spomínané v predchádzajúcej kapitole, v časti \ref{kap2:2.2:2.2.3:Tree}. Toto kódovanie je ale nutné upraviť tak, aby bralo do úvahy zvyšok definície rozhodovacieho stromu spolu s kritériami delenia zadefinovanými v oddieli \ref{kap1:2.5:DTSplitCriterias}.

Vnútorný uzol $u$ rozhodovacieho stromu je kódovaný trojicou $(a,z,h)$, kde $a$ predstavuje atribút delenia, $z$ porovnávací operátor použitý pri delení a $v$ hodnotu, v ktorej delíme dátovú množinu. Vieme, že typ atribútu môže byť numerický alebo kategoriálny. Podľa typu atribútu $a$ v trojici $(a,z,h)$ sú $z$ a $h$ definované takto:
\begin{enumerate}
\item V prípade numerického atribútu $a$ je $z$ porovnávací operátor a $h$ reálna hodnota.
\item V prípade kategoriálneho atribútu $a$ sú $z$ a $h$ nedefinované. Každá hrana vychádzajúca z vrcholu zodpovedá jednej možnej hodnote atribútu $a$.
\end{enumerate}
Počet hrán vychádzajúcich z vrcholu $u$ sa riadi typom atribútu $a$.
\begin{enumerate}
\item V prípade numerického atribútu sú hrany dve.
\item V prípade kategoriálneho atribútu je počet hrán rovný veľkosti definičného oboru atribútu $a$.
\end{enumerate}
Listový uzol $l$ rozhodovacieho stromu obsahuje iba hodnotu $c$, ktorá predstavuje klasifikáciu.
Takto kódovaný rozhodovací strom je na Obrázku \ref{fig:DTEncode}.

\begin{figure}[h]
\centering
\centerline{\mbox{\includegraphics{../img/kap3/dtgen.pdf}}}
\caption{Príklad kódovania rozhodovacieho stromu v genetickom algoritme na tvorbu rozhodovacích stromov. Nech hodnoty výstupného atribútu sú z množiny $C=\{c_0,c_1\}$. Na obrázku sú tri vnútorné uzly s atribútami $x,y,z$. Atribúty $x$ a $y$ sú numerické a $z$ je kategoriálny. Nedefinovaná hodnota operátora a nedefinovaná hodnota atribútu (u kategoriálneho atribútu) sú označené *. Hrany predstavujú rozhodnutia a počet hrán závisí na type atribútu uzlu. Listové vrcholy obsahujú hodnoty z množiny $C$ (klasifikáciu).}\label{fig:DTEncode}
\end{figure}

\section{Fitness funkcie}\label{kap3:3.3:Fitness}
Fitness funkcie v genetickom algoritme určujú, ako často budú jedinci vyberaní k reprodukcii. Tieto funkcie volíme podľa toho, čo chceme v rozhodovacom strome optimalizovať. Voľba vhodnej fitness funkcie v genetickom algoritme ovplyvňuje
\begin{itemize}
\item prediktívne vlastnosti stromu, kam zaraďujeme trénovaciu, testovaciu a generalizačnú chybu, 
\item zložitosť stromu, meranú napríklad výškou stromu, počtom listov alebo počtom vrcholov.
\end{itemize} 

Na optimalizovanie prediktívnych vlastností stromu musíme voliť také fitness funkcie, ktoré s predikciou súvisia. Tieto fitness funkcie väčšinou počítajú určitú metriku z matice chybovosti. Na zlepšenie prediktívnych vlastností nie je každá metrika vhodná. Kvôli tomu je lepšie sa zamerať iba na niektoré z nich.
V časti \ref{kap1:2.6:2.6.2:Alternatives} sme definovali alternatívne metriky založené na matici chybovosti, ktoré môžeme použiť pri ohodnocovaní rozhodovacích stromov. Spomínané kriteriá citlivosti, precíznosti, špecificity a f-metriky sú vhodnými kandidátmi fitness funkcií. Na zlepšenie generalizačnej chyby by mal byť výpočet metrík vykonaný na validačnej množine.

Z časti \ref{kap2:2.2:2.2.3:Tree} vieme, že stromy v genetických algoritmoch majú tendenciu narastať a je nutné ich veľkosť obmedziť. To sa dá dosiahnuť voľbou fitness funkcií, ktoré ohodnocujú zložitosť stromu. Takéto fitness funkcie využívajú metriky zadefinované v časti \ref{kap1:2.6:2.6.3:TreeLook}. K existujúcim funkciám môžeme definovať nové fitness funkcie, ktoré budeme nazývať \emph{relatívne}. Takéto funkcie sú počítané vzhľadom k nejakej hodnote \emph{x}, v ktorej zároveň dosahujú maximum. Na Obrázku \ref{fig:normalvsrelative} uvádzame grafy dvoch fitness funkcií (normálnej a relatívnej) pre veľkosť stromov.

Vzorec na výpočet relatívnej fitness funkcie musí pracovať s hodnotou $y$, od ktorej túto fitness funkciu chceme dopočítať. Nech $x$ je skutočná hodnota fitness funkcie. Relatívnu fitness funkciu môžeme definovať jedným zo vzorcov
\begin{itemize}
\item $\dfrac{1}{2^{|x-y|}}$ 
\item $\dfrac{1}{2^{(x-y)^2}}$
\item gausovskou funkciou $\mbox{exp}\left(-\dfrac{(x-y)^2}{2z^2}\right)$
\end{itemize}

\begin{figure}[h]
\centering
\centerline{\mbox{\includegraphics[width=400pt]{../img/kap3/normvsrel.pdf}}}
\caption{Príklad grafov normálnej a relatívnej fitness funkcie hodnotiace veľkosť stromov. Na ľavo je normálna fitness funkcia, ktorá dosahuje maximum pre najmenšiu veľkosť stromu a vpravo je relatívna fitness funkcia s maximom v hodnote \emph{x}. Vzorec ľavej fitness funkcie je $\left(\dfrac{1}{x}\right)$ a vzorec pravej fitness funkcie je $\left(\dfrac{1}{2^{x-y}}\right)$.}\label{fig:normalvsrelative}
\end{figure}

Voľba správnych fitness funkcií je obtiažnou úlohou. Optimálne by mali zvolené funkcie zmenšovať generalizačnú chybu a zbytočne nekomplikovať stromy. Tým sa úloha stáva multi-kriteriálnym problémom a výpočet kvality jedinca sa robí jednou z metód spomínaných v časti \ref{kap2:2.3:Fitnesses}.
\section{Operátory na tvorbu rozhodovacích stromov}\label{kap3:3.4:Operators}
V oddieli \ref{kap2:2.5:Operators} sme zadefinovali genetické operátory a uviedli ich dôležitosť v genetických algoritmoch. V prípade genetickej tvorby rozhodovacích stromov sú operátory zodpovedné za vytváranie nových rozhodovacích stromov. Tieto operátory by mali pracovať v súlade s dátovými množinami. To znamená nevykonávať náhodné zmeny, ale iba také, ktoré nejako súvisia s dátami.

Z predchádzajúcich častí vieme, že v genetických algoritmoch obvykle najprv vykonávame kríženie a potom mutáciu. Taktiež vieme, že každý operátor aplikujeme iba s určitou pravdepodobnosťou a že toto nastavenie je neľahkou úlohou. Jednou z možností je inšpirovať sa oddielom \ref{kap2:2.6:Parameters}.

V časti \ref{kap3:3.4:3.4.1:Crossover} pripomenieme kríženie podstromov. Potom v časti \ref{kap3:3.4:3.4.2:Mutation} predvedieme dva nové druhy mutácií.

\subsection{Kríženie}\label{kap3:3.4:3.4.1:Crossover}
V časti \ref{kap2:2.5:2.5.1:Crossover} sme definovali operátor kríženia podstromov. V prípade stromov je toto jediné kríženie, ktoré dáva zmysel, a preto ho budeme používať aj pri genetickej tvorbe rozhodovacích stromov.

Toto kríženie ale musíme upraviť tak, aby v prípade kategoriálneho atribútu kontrolovalo prekrížené stromy. Krížený podstrom môže obsahovať vrchol, ktorého atribút sa už vyskytuje na ceste do koreňa a hodnoty atribútu sú nekonzistentné. V takomto prípade ho nahradím správnym potomkom (s konzistentnou hodnotou atribútu). Po zmene znova vykonám kontrolu a opakujem tento postup, dokým nie sú tieto atribúty konzistentné.

\subsection{Mutácie}\label{kap3:3.4:3.4.2:Mutation}
V časti \ref{kap2:2.5:2.5.2:Mutation} sme uviedli rôzne druhy mutácií jedincov kódovaných stromami. Po menších úpravách sa dajú tieto mutácie použiť v genetickom algoritme na tvorbu rozhodovacích stromov. Mutácie je väčšinou nutné upraviť tak, aby brali do úvahy kódovanie uvedené v časti \ref{kap3:3.2:Encoding}. Jednotlivé mutácie by tiež mali vykonávať zmeny, ktoré sú v súlade s danou dátovou množinou $D$ (zabránenie náhodným zmenám).

Jednou zo spomínaných mutácií, ktorú použijeme v genetickom algoritme, je zmena náhodného vnútorného uzlu na list. Táto mutácia teda odstráni všetky odchádzajúce hrany zo zvoleného vrcholu a nastaví hodnotu klasifikácie.

Nech $S$ je strom, na ktorom vykonávame mutáciu, $u$ je \emph{vnútorný} uzol stromu $S$. Ďalej nech $D_u$ je podmnožina množiny $D$ získaná delením $D$ podľa vrcholov (konkrétne atribútov a hodnôt vo vrchole) na ceste z koreňa do vrcholu $u$ a $\max_c u$ je najčastejšia hodnota výstupného atribútu v $D_u$. Potom jednotlivé kroky mutácie sú
\begin{enumerate}
\item zvoľ vnútorný uzol $u$,
\item nahraď $u$ listom s klasifikáciou $\max_c u$.
\end{enumerate}
Príklad takejto mutácie nahradzujúcej podstrom listom je znázornený na Obrázku \ref{fig:mutnodetoleaf}.

\begin{figure}[h]
\centering
\centerline{\mbox{\includegraphics[width=400pt]{../img/kap3/mutnodetoleaf.pdf}}}
\caption{Príklad mutácie rozhodovacieho stromu, ktorá mení vnútorný uzol na list. Zvolený vnútorný uzol $u$ (nech $\max_c u = c_0$) je označený čiarkovaným obdĺžnikom. Na obrázku vpravo je vnútorný uzol nahradení listom s hodnotou $c_0$ (klasifikáciou).}\label{fig:mutnodetoleaf}
\end{figure}

Druhou použitou mutáciou je priradenie nového podstromu do nejakého uzlu stromu. Podstromy budú generované z podmnožín dátovej množiny $D$ a uložené v predom definovanej množine podstromov. Počet podmnožín bude rovný počtu vytváraných podstromov. Tieto podmnožiny môžeme vytvárať:
\begin{enumerate}
\item delením dátovej množiny $D$ na toľko častí, koľko podstromov budeme vytvárať,
\item výberom s opakovaním z dátovej množiny $D$.
\end{enumerate}
Nové podstromy na priradenie by nemali byť náhodné a ani zbytočne veľké. Mutácia preto používa rozhodovacie korene (definícia \ref{kap1:2.3:2.3.2:stumpDT}).

Nech $K=\{K_1,\ldots,K_r\}$ je množina podstromov (rozhodovacích koreňov). Ďalej nech $S$ je strom, na ktorom vykonávame mutáciu, $u$ ľubovoľný uzol stromu $S$. Potom jednotlivé kroky mutácie sú 
\begin{enumerate}
\item vyber náhodne podstrom $K_i$ z $K$,
\item vyber náhodne vrchol $u$ stromu $S$,
\item nahraď $u$ podstromom $K_i$.
\end{enumerate}
Príklad mutácie nahradzujúcej podstrom rozhodovacím koreňom je na Obrázku \ref{fig:mutdecisionstump}.

Podobne ako pri krížení je treba túto mutáciu upraviť, aby v prípade kategoriálneho atribútu kontrolovala, či sa atribút vrcholu rozhodovacieho koreňa (ten je iba jeden) už nevyskytuje na ceste do koreňa stromu a hodnoty atribútu by boli nekonzistentné. V takom prípade nahradíme vrchol jeho potomkom  (čo je určite list) pre správny atribút.

\begin{figure}[h]
\centering
\centerline{\mbox{\includegraphics[width=400pt]{../img/kap3/mutdecisionstump.pdf}}}
\caption{Príklad mutácie rozhodovacieho stromu pomocou rozhodovacích koreňov. Na obrázku hore je množina vygenerovaných rozhodovacích koreňov $K$. Na obrázku dole je samotná mutácia stromu. Uzol určený na mutáciu je označený čiarkovaným štvorcom. Na obrázku (dole vpravo) je zvolený vrchol $u$ (list) nahradený rozhodovacím koreňom $K_1$. }\label{fig:mutdecisionstump}
\end{figure}

% Ukázka použití některých konstrukcí LateXu (odkomentujte, chcete-li)
% %%% Ukázka použití některých konstrukcí LaTeXu

\subsection{Ukázka \LaTeX{}u}
\label{ssec:ukazka}

V~této krátké části ukážeme použití několika základních konstrukcí \LaTeX{}u,
které by se vám mohly při psaní práce hodit.

Třeba odrážky:

\begin{itemize}
\item Logo Matfyzu vidíme na obrázku~\ref{fig:mff}.
\item Tato subsekce má číslo~\ref{ssec:ukazka}.
\item Odkaz na literaturu~\cite{lamport94}.
\end{itemize}

Druhy pomlček:
červeno-černý (krátká),
strana 16--22 (střední),
$45-44$ (minus),
a~toto je --- jak se asi dalo čekat --- vložená věta ohraničená dlouhými pomlčkami.
(Všimněte si, že jsme za \verb|a| napsali vlnovku místo mezery: to aby se
tam nemohl rozdělit řádek.)

% Makro na české uvozovky (novější verze LaTeXu ho už mají zabudované)
\newcommand{\uv}[1]{\quotedblbase #1\textquotedblleft}
\uv{České uvozovky.}

\newtheorem{theorem}{Věta}
\newtheorem*{define}{Definice}	% Definice nečíslujeme, proto "*"

\begin{define}
{\sl Strom} je souvislý graf bez kružnic.
\end{define}

\begin{theorem}
Tato věta neplatí.
\end{theorem}

\begin{proof}
Neplatné věty nemají důkaz.
\end{proof}

\begin{figure}
	\centering
	\includegraphics[width=30mm]{../img/logo.eps}
	\caption{Logo MFF UK}
	\label{fig:mff}
\end{figure}


\chapter{Záver}\label{kap:fin}
\renewcommand{\figurename}{Obrázok}
Všetky ciele práce sa podarilo bezozvyšku splniť. Genetický algoritmus na vytváranie rozhodovacích stromov, ktorý sme navrhli a naprogramovali v tejto práci dokáže skonštruovať rozhodovacie stromy. Kvalitou sú tieto stromy porovnateľné so stromami vytvorenými algoritmom C4.5 a zároveň sú od nich výrazne menšie.

Implementovaný algoritmus obsahuje značné množstvo nastaviteľných parametrov, ktoré je možné meniť pri každom spustení programu. Na druhú stranu sme sa v práci hlbšie nezaoberali hľadaním optimálneho nastavenia parametrov pre konkrétnu úlohu, pretože veľkosť priestoru parametrov je príliš veľký. Preskúmanie vhodných parametrov by vyžadovalo rozsiahle experimenty skúmajúce dopad rôznych nastavení na výsledné stromy.

V práci sme porovnali nový genetický algoritmus tvorby rozhodovacích stromov s indukčným algoritmom C4.5 a dopadlo to dobre.
Pretože sme od algoritmu vyžadovali, aby vytváral menšie stromy, tak sme ich získali za cenu nepatrného zhoršenia generalizačných schopností vytvorených stromov a pravdaže za cenu dlhšieho čakania na výsledky.

Aplikácia \verb|GenDTLib| je spustiteľná a použiteľná z nástroja Weka. Implementácia je spoľahlivá a rýchlosť tvorby modelu je postačujúca, keď berieme do úvahy, že používame genetické algoritmy. Na implementovanie sme použili jazyk Java, ktorá poskytuje nezávislosť aplikácie od operačného systému. S aplikáciou sme navrhli aj aplikačné rozhranie, pri vytváraní ktorého sme sa snažili o to, aby bolo zrozumiteľné a v prípade ďalšieho rozšírenia sa sním pracovalo čo najľahšie. Zároveň sme museli zohľadniť to, že aplikácia má byť kompatibilná s nástrojom Weka. Parametre genetického algoritmu je možné nastaviť z konfiguračného súboru a existujúce komponenty je taktiež možné rozšíriť o nové pomocou zásuvných modulov.

Túto aplikáciu je možné spustiť na bežne dostupných počítačoch, na ktorých je spustiteľný aj nástroj Weka. 
Voľba vytvoriť aplikáciu ako zásuvný modul do nástroja Weka uľahčila prácu pri testovaní a umožňuje použiť klasifikátor z grafického rozhrania. Skonštruované stromy môžeme navyše vďaka tomu zobraziť zo vstavaného vizualizačného nástroja (obr. \ref{fig:wekatree}).

\begin{figure}[h]
\centering
\centerline{\mbox{\includegraphics[width=300pt]{../img/zaver/visualize.pdf}}}
\caption{Obrazovka z nástroja Weka, na ktorej je vizualizačný nástroj na zobrazovanie rozhodovacích stromov. Štruktúra stromov je podľa Definície \ref{kap1:2.3:2.3.2:DT}}\label{fig:wekatree}
\end{figure}

\subsection{Možné rozšírenia práce}
Aplikácia \verb|GenDTLib| je v súčasnej dobe plnohodnotný klasifikátor v nástroji Weka, ktorý vytvára rozhodovacie stromy pomocou genetických algoritmov. Možnosť nastaviť parametre genetického algoritmu v konfiguračnom súbore umožňuje užívateľovi meniť a experimentovať s existujúcimi komponentami a ich nastaveniami. Aplikáciu je ale možné rozšíriť aj o nové komponenty, ktoré sme v práci neuvažovali. Pri ďalšom vývoji alebo v budúcich verziách aplikácie by bolo vhodné pridať niektoré funkcie:

\begin{itemize}
\item \textbf{Nový kontajner pre populáciu}. V pokročilých genetických algoritmoch sa problém správnej voľby veľkosti populácie rieši dynamicky sa meniacou veľkosťou populácie (na začiatku väčšia a postupne sa zmenšuje). Aplikáciu je možné rozšíriť o nové druhy populácií a preto by bolo v budúcnosti vhodné takýto druh dynamickej populácie implementovať.
\item \textbf{Lepšia paralelizácia}. Aplikácia v aktuálnej verzii umožňuje nastaviť počet vlákien, ktoré počítajú fitness funkcie a generujú počiatočnú populáciu, no do budúcna by bolo možno vhodné paralelizovať aj niektoré komplikovanejšie operátory. Ďalšou možnosťou je implementovať niektorú z techník na paralelizovanie genetických algoritmov (napr. ostrovný model, hybridný paralelný genetický algoritmus).
\item \textbf{Vylepšiť zastavovacie kritéria}. Najpoužívanejším kritériom zastavenia je pri genetických algoritmoch počet generácií. V aktuálnej verzii je algoritmus zastavený po pevne danom počte generácií. Kritériá zastavenia môžu byť ale vhodnejšie. Pri ďalšom vývoji aplikácie by bolo vhodné vyskúšať a otestovať kvalitu riešení a rýchlosť genetických algoritmov s upravenými kritériami zastavenia. Takáto zmena by takisto mohla urýchliť celkový beh genetických algoritmov.
\item \textbf{Pareto fitness}. Pri niektorých úlohách riešených genetickými algoritmami sa používa pareto fitness. V aktuálnej verzii aplikácie je tento druh fitness nepodporovaný, a preto by bolo vhodné pareto fitness v budúcich verziách pridať a otestovať.
\end{itemize}

%%% Seznam použité literatury
%%% Seznam použité literatury je zpracován podle platných standardů. Povinnou citační
%%% normou pro diplomovou práci je ISO 690. Jména časopisů lze uvádět zkráceně, ale jen
%%% v kodifikované podobě. Všechny použité zdroje a prameny musí být řádně citoványy

\def\bibname{Zoznam použitej literatúry}
\addcontentsline{toc}{chapter}{\bibname}
\printbibliography
% http://en.wikipedia.org/wiki/Data_mining %



%%% Tabulky v diplomové práci, existují-li.
\chapwithtoc{Zoznam tabuliek}

%%% Použité zkratky v diplomové práci, existují-li, včetně jejich vysvětlení.
\chapwithtoc{Zoznam použitých skratiek}
\textbf{KDD} - Knowledge Discovery in Databases \\
\textbf{DM} - Data mining / Dobývanie znalostí \\
\textbf{CRISP-DM} - CRoss-Industry Standard Process for Data Mining \\
\textbf{SVM} - Support vector machine \\
\textbf{AID} - Automatic Interaction Detection \\
\textbf{THAID} - THeta Automatic Interaction Detection \\
\textbf{CHAID} - CHi-squared Automatic Interaction Detection \\
\textbf{ID3} - Iterative Dichotomiser 3 \\
\textbf{CART} - Classification And Regression Tree \\
\textbf{TDIDT} - Top down induction of decision trees \\
\textbf{SCLGSI} - Symptom Checklist-90-R Global Severity Index
\textbf{GA} - Genetické algoritmy 

%%% Přílohy k diplomové práci, existují-li (různé dodatky jako výpisy programů,
%%% diagramy apod.). Každá příloha musí být alespoň jednou odkazována z vlastního
%%% textu práce. Přílohy se číslují.
\chapwithtoc{Prílohy}

\openright
\end{document}
